% !TEX program = xelatex
% =====================================================================
%  BÁO CÁO ĐỒ ÁN TRÍ TUỆ NHÂN TẠO, HỌC PHẦN AIP202
%  Khoa Công nghệ Thông tin, Trường Đại học Kiến trúc Đà Nẵng
%
%  Build:  latexmk -xelatex do-an.tex
%  Overleaf: Menu > Compiler > XeLaTeX, rồi bấm Recompile.
% =====================================================================
% \documentclass phải nằm ở file này, không được chuyển vào settings.tex,
% vì Overleaf dò file chính bằng cách tìm file có chứa \documentclass.
\documentclass[12pt, a4paper, oneside]{book}
% !TEX root = ../do-an.tex
% =====================================================================
%  styles/settings.tex
%  Các thiết lập chung, nạp từ file chính do-an.tex.
%  SINH VIÊN KHÔNG SỬA FILE NÀY.
%
%  Lưu ý: lệnh khai báo lớp tài liệu nằm ở file do-an.tex chứ không nằm
%  ở đây. Lý do: Overleaf dò file chính bằng cách tìm chuỗi khai báo lớp
%  tài liệu trong các file. Nếu chuỗi đó xuất hiện ở file này, Overleaf
%  sẽ biên dịch nhầm file này và báo lỗi "no legal end found", vì file
%  này chỉ chứa thiết lập chứ không chứa phần thân tài liệu.
% =====================================================================

% Toàn bộ nhận diện DAU: phông, màu, bìa, hộp trình bày, header
\usepackage{styles/dau}
\dauLayout

% Giãn dòng và khoảng cách đoạn theo quy cách đồ án
\linespread{1.25}
\usepackage[skip=8pt plus 1pt, indent=15pt]{parskip}

% Ghi chú: mẫu gốc dùng gói atbegshi để vứt trang đầu, vì trang bìa của
% mẫu gốc dựng bằng minipage nên sinh ra một trang trắng. Trang bìa ở
% đây dựng bằng môi trường titlepage nên không có trang trắng, và nếu
% giữ thủ thuật đó thì chính trang bìa sẽ bị vứt mất.


% =====================================================================
%  THÔNG TIN ĐỒ ÁN — thay các trường trong ngoặc vuông trước khi nộp
% =====================================================================
% Loại báo cáo
\crname{BÁO CÁO ĐỒ ÁN TRÍ TUỆ NHÂN TẠO}
% Học phần
\cmonhoc{Học phần AIP202 \textperiodcentered{} 2 tín chỉ}
% Tên đề tài. Dùng \\ để ngắt dòng cho cân đối trên bìa.
\ctname{XÂY DỰNG NỀN TẢNG AI EXAMHUB DAU\\HỖ TRỢ TẠO ĐỀ VÀ TỔ CHỨC THI TRỰC TUYẾN}
% Họ tên sinh viên và mã số sinh viên
\cstuname{Sinh viên 1 \textperiodcentered{} Sinh viên 2}
\cmssv{Đã ẩn danh}
\cstudentrows{Sinh viên 1 & Đã ẩn danh\\ Sinh viên 2 & Đã ẩn danh\\}
% Lớp
\clopsv{23CT1}
% Ngành
\csCouncil{Công nghệ thông tin}
% Giảng viên hướng dẫn
\csSupervise{Đã ẩn danh}
% Năm học
\cttime{2026--2027}
% Khoa quản lý
\csdeptname{KHOA CÔNG NGHỆ THÔNG TIN}
% Nơi thực hiện
\csplace{Đà Nẵng}

\begin{document}

\setcounter{tocdepth}{2}
\setcounter{secnumdepth}{3}

% Bìa trang trọng có khung kép. Muốn dùng bìa hiện đại theo bộ giáo
% trình của khoa thì thay dòng dưới bằng \coverpagehiendai
\coverpage

% ---------------------------------------------------------------------
%  PHẦN ĐẦU: đánh số trang bằng số La Mã thường
% ---------------------------------------------------------------------
\frontmatter
\dauPageStyle

% ---------------------------------------------------------------------
%  LỜI CAM ĐOAN
% ---------------------------------------------------------------------
\begin{declaration}

Tôi xin cam đoan báo cáo này được thực hiện trong quá trình xây dựng và
khảo sát hệ thống AI ExamHub DAU. Nội dung mô tả thiết kế, chức năng và kết
quả kiểm thử được đối chiếu với mã nguồn tại thời điểm hoàn thành báo cáo.
Các số liệu được nêu có căn cứ từ kết quả chạy thử; báo cáo không xem các
giá trị cấu hình hoặc điểm tin cậy do chương trình trả về là kết quả đánh
giá độ chính xác của mô hình.

Mọi số liệu, hình ảnh, mã nguồn và kết quả của tác giả khác được sử dụng
trong đồ án đều được trích dẫn nguồn gốc rõ ràng trong phần tài liệu tham
khảo. Toàn bộ phần mềm và thư viện được dùng đều có giấy phép cho phép sử
dụng. Những phần có sử dụng công cụ trí tuệ nhân tạo hỗ trợ đã được nêu rõ
trong báo cáo.

Nếu phát hiện có bất kỳ sự gian lận nào, tôi xin chịu hoàn toàn trách nhiệm
trước Khoa và Nhà trường.

\vspace{1.4em}
\hfill\okyten{\dauSinhVien}{Sinh viên thực hiện}

\end{declaration}

% ---------------------------------------------------------------------
%  LỜI CẢM ƠN
% ---------------------------------------------------------------------
\begin{acknowledgments}

Em xin trân trọng cảm ơn quý thầy cô Khoa Công nghệ Thông tin, Trường Đại
học Kiến trúc Đà Nẵng đã trang bị kiến thức nền tảng và tạo điều kiện để em
thực hiện đồ án. Em cũng xin cảm ơn giảng viên hướng dẫn đã dành thời gian
góp ý, giúp em định hướng việc phân tích yêu cầu, triển khai hệ thống và
hoàn thiện báo cáo. Cuối cùng, em xin cảm ơn gia đình và bạn bè đã động
viên em trong quá trình học tập và thực hiện đồ án.

\vspace{1.4em}
\hfill\okyten{\dauSinhVien}{Sinh viên thực hiện}

\end{acknowledgments}

% ---------------------------------------------------------------------
%  TÓM TẮT
% ---------------------------------------------------------------------
\begin{abstract}

AI ExamHub DAU là nền tảng web hỗ trợ giảng viên quản lý ngân hàng câu hỏi,
tạo đề từ tài liệu học tập và tổ chức thi trực tuyến cho sinh viên. Hệ thống
được xây dựng theo kiến trúc web gồm giao diện React và TypeScript, API
Node.js/Express và cơ sở dữ liệu PostgreSQL. Tài liệu PDF, DOC/DOCX và TXT
được tiếp nhận để trích xuất văn bản; khi cấu hình khóa dịch vụ, hệ thống có
thể yêu cầu mô hình ngôn ngữ tạo câu hỏi, còn khi không có dịch vụ thì sử
dụng bộ sinh dự phòng dựa trên nội dung nguồn. Giảng viên có thể chọn loại
câu hỏi, độ khó, mức độ tư duy, rà soát căn cứ trích dẫn, lập quy trình duyệt
đề và tạo phòng thi. Sinh viên tham gia bằng mã phòng, lưu và tiếp tục bài
làm; hệ thống tự chấm câu trắc nghiệm và cho phép giảng viên chấm tự luận.
Quá trình kiểm chứng gồm 35 bài kiểm thử API, đều đạt, cùng một lượt build
TypeScript/Vite thành công. Đây là bằng chứng kiểm thử chức năng và build,
không phải phép đo độ chính xác của AI: đồ án chưa có tập dữ liệu chuẩn hay
benchmark độc lập để định lượng chất lượng câu hỏi, OCR hoặc hiệu năng. Kết
quả là một nguyên mẫu tích hợp được quy trình tạo đề và thi, đồng thời làm
rõ các phần cần tiếp tục đánh giá trước khi triển khai thực tế.

\vspace{0.8em}
\noindent\textbf{Từ khoá:} tạo câu hỏi, xử lý tài liệu, thi trực tuyến,
React, Express, PostgreSQL.
\end{abstract}


\tableofcontents
\listoffigures
\listoftables

% ---------------------------------------------------------------------
%  THÂN BÁO CÁO
% ---------------------------------------------------------------------
\mainmatter
\dauPageStyle

% ---------------------------------------------------------------------
%  CHƯƠNG 1 — MỞ ĐẦU
% ---------------------------------------------------------------------
\chapter{Mở đầu và cơ sở lý thuyết}
\label{ch:modau}

\begin{muctieu}
Chương này nêu bối cảnh, mục tiêu và phạm vi đồ án, sau đó trình bày cơ sở
lý thuyết cùng công nghệ của nền tảng AI ExamHub DAU.
\end{muctieu}

\section{Lý do chọn đề tài}

Trong quy trình đào tạo, giảng viên cần biên soạn câu hỏi, tạo nhiều đề,
quản lý lịch thi, theo dõi bài làm và tổng hợp kết quả. Khi các công việc
này thực hiện rời rạc, việc tái sử dụng học liệu và theo dõi trạng thái
kiểm tra trở nên khó khăn. AI ExamHub DAU được xây dựng như một nền tảng
thống nhất cho ngân hàng câu hỏi, hỗ trợ sinh câu hỏi từ tài liệu, quy trình
duyệt đề và tổ chức lớp học, phòng thi trực tuyến.

Đề tài phù hợp với học phần trí tuệ nhân tạo vì kết hợp một ứng dụng thực
tiễn với thành phần xử lý ngôn ngữ: đọc nội dung học liệu, đề xuất câu hỏi
và liên kết câu hỏi với bằng chứng trong nguồn. Hệ thống cũng có phương án
dự phòng dựa trên quy tắc khi không cấu hình dịch vụ mô hình ngôn ngữ. Việc
khảo sát đồng thời phần AI và quy trình web giúp làm rõ giới hạn của một
nguyên mẫu tích hợp, thay vì đồng nhất việc kết nối API với một mô hình đã
được huấn luyện và đánh giá độc lập.

\section{Mục tiêu của đồ án}

Đồ án hướng tới các mục tiêu sau:
\begin{itemize}[leftmargin=1.6em, itemsep=2pt]
  \item Xây dựng ứng dụng web gồm giao diện người dùng, API backend và cơ sở
        dữ liệu quan hệ.
  \item Hỗ trợ tải tài liệu học tập, trích xuất văn bản và đề xuất câu hỏi
        trắc nghiệm/tự luận có thông tin căn cứ.
  \item Hỗ trợ quản lý câu hỏi, tạo và duyệt đề thi, tổ chức phòng thi trực
        tuyến và lớp học theo vai trò.
  \item Kiểm tra các luồng chính bằng kiểm thử API và xác nhận giao diện có
        thể build/type-check.
  \item Nêu rõ các giới hạn của phần AI, dữ liệu và kiểm thử, không suy diễn
        kết quả chức năng thành độ chính xác mô hình.
\end{itemize}

\section{Câu hỏi nghiên cứu và tiêu chí thành công}

Phạm vi đồ án được định hướng bởi ba câu hỏi. Thứ nhất, làm sao kết nối việc
tạo đề, duyệt đề và tổ chức kiểm tra vào cùng một quy trình phần mềm có phân
quyền? Thứ hai, ở mức một nguyên mẫu, làm thế nào hỗ trợ khai thác tài liệu
học tập mà vẫn giữ được khả năng giảng viên rà soát nguồn? Thứ ba, những
bằng chứng nào là đủ để kết luận về độ hoàn thiện phần mềm, và những câu hỏi
nào cần một nghiên cứu đánh giá AI riêng?

Tiêu chí thành công của phần mềm là các chức năng trọng tâm có API/UI tương
ứng, quyền truy cập được phân theo vai trò, dữ liệu chính được lưu trong
PostgreSQL, các luồng nghiệp vụ có kiểm thử tự động và frontend build
thành công. Với thành phần AI, tiêu chí của nguyên mẫu chỉ là tạo được đầu
ra có cấu trúc hoặc nhánh dự phòng và cung cấp thông tin nguồn để kiểm tra.
Độ đúng, tính sư phạm, độ khó và hiệu quả học tập không được xem là đã xác
nhận cho đến khi có bộ chuẩn đánh giá độc lập.

\section{Phương pháp thực hiện}

Đồ án sử dụng phương pháp phát triển theo hướng tích hợp và kiểm chứng
nghiệp vụ. Trước hết khảo sát cấu trúc mã nguồn và schema, xác định các vai
trò cùng luồng sử dụng; tiếp đó mô tả kiến trúc client--server, mô hình dữ
liệu và hợp đồng API. Sau khi hiện thực frontend, backend và các migration,
các luồng được xác minh bằng kiểm thử API hiện có và lệnh build của
TypeScript/Vite. Cuối cùng, kết quả được phân loại thành bằng chứng chức
năng, bằng chứng build và các đánh giá chưa có dữ liệu hỗ trợ.

Phương pháp này phù hợp với mục tiêu tạo nguyên mẫu tích hợp trong phạm vi
đồ án, nhưng không thể thay thế thực nghiệm có người dùng, đo tải, kiểm thử
bảo mật độc lập hoặc đánh giá chất lượng nội dung AI bởi chuyên gia. Phạm
vi và loại bằng chứng được tóm tắt trong Bảng~\ref{tab:phuongphap}.

\begin{table}[H]
\centering
\caption{Phương pháp và bằng chứng tương ứng với câu hỏi của đồ án.}
\label{tab:phuongphap}
\begin{tabularx}{\textwidth}{>{\bfseries\color{daunavy}}p{0.25\textwidth} X X}
\toprule
\rowcolor{gray100}Câu hỏi & Cách khảo sát & Bằng chứng có/không có \\
\midrule
Tích hợp nghiệp vụ & Đọc frontend, route API, middleware và schema; chạy luồng kiểm thử. & Mã nguồn, 35 ca API đạt và build frontend thành công. \\
Hỗ trợ khai thác tài liệu & Đọc xử lý trích xuất, cấu hình mô hình, kiểm tra đầu ra và fallback. & Có kiểm thử chức năng; chưa có chấm điểm câu hỏi bởi chuyên gia. \\
Mức độ chất lượng AI & Xác định tập chuẩn và độ đo cần thiết. & Chưa có tập gán nhãn; không báo cáo accuracy/F1/độ đúng nội dung. \\
\bottomrule
\end{tabularx}
\end{table}

\section{Đóng góp và giới hạn đóng góp}

Đóng góp chính là tích hợp được quy trình tạo câu hỏi có căn cứ, chọn ngân
hàng để lập đề, duyệt và phát hành, tổ chức phòng thi, lưu tiến độ của
sinh viên và ghi nhận kết quả trên cùng nền tảng. Về kỹ thuật, hệ thống
minh hoạ cách phối hợp giao diện React/TypeScript, REST API Express và
PostgreSQL, có Docker Compose cho môi trường phát triển. Thành phần AI được
đóng góp ở mức tích hợp suy luận dịch vụ ngoài khi cấu hình cùng một nhánh
sinh dự phòng; đồ án không tuyên bố đề xuất kiến trúc mô hình mới, huấn
luyện mô hình riêng hay tạo ra một corpus học thuật mới.

\section{Phạm vi, đối tượng và giả định}

Đối tượng sử dụng gồm quản trị viên, giảng viên và sinh viên. Phạm vi triển
khai khảo sát trong báo cáo tập trung vào ứng dụng web, API Node.js/Express
và PostgreSQL. Luồng tạo câu hỏi nhận nội dung tài liệu PDF, DOC/DOCX hoặc
TXT; khả năng trích xuất phụ thuộc vào định dạng và chất lượng văn bản/OCR.
Mô hình ngôn ngữ bên ngoài chỉ được sử dụng khi nhà triển khai cung cấp cấu
hình cần thiết; nếu không, hệ thống dùng bộ sinh dự phòng dựa trên văn bản.

Đồ án không xây dựng hay huấn luyện mô hình nền tảng riêng, không có bộ dữ
liệu gán nhãn chuẩn để đo độ chính xác sinh câu hỏi, không triển khai các
thư mục AI riêng như những dịch vụ độc lập, và chưa đánh giá tải hoặc bảo
mật ở quy mô sản xuất. Chấm điểm tự động hiện tập trung vào câu trắc nghiệm;
bài tự luận cần giảng viên chấm theo nội dung và tiêu chí.

Các giả định gồm: người tạo đề có quyền sử dụng tài liệu được tải lên; nhà
triển khai tự cấu hình dịch vụ AI nếu muốn sử dụng nhánh mô hình ngôn ngữ;
giảng viên xem xét câu hỏi/đáp án trước khi phát hành; cơ sở dữ liệu được
khởi tạo theo schema/migration tương ứng. Những giả định này cần được rà
soát khi đưa hệ thống vào sử dụng tại một cơ sở đào tạo thật.

\section{Bố cục báo cáo}

Phần còn lại của báo cáo gồm ba chương. Chương~1 trình bày cơ sở lý thuyết
và công nghệ liên quan; Chương~2 phân tích yêu cầu, thiết kế dữ liệu và
kiến trúc, đồng thời mô tả phần cài đặt; Chương~3 tổng hợp kết quả, bằng
chứng kiểm thử, hạn chế và hướng phát triển. Sau phần tài liệu tham khảo là
phụ lục hướng dẫn chạy hệ thống, từ điển dữ liệu và danh mục API.

% ---------------------------------------------------------------------
%  NỘI DUNG CHƯƠNG 1 — CƠ SỞ LÝ THUYẾT VÀ CÔNG NGHỆ
% ---------------------------------------------------------------------
\section{Cơ sở lý thuyết và công nghệ}
\label{sec:coso}

\begin{muctieu}
Phần này giới thiệu quy trình sinh câu hỏi từ tài liệu, các thành phần xử
lý dữ liệu và công nghệ thực tế trong hệ thống.
\end{muctieu}

\subsection{Kiến thức nền tảng}

Trong phạm vi đồ án, AI được dùng để hỗ trợ suy luận và sinh nội dung từ
học liệu qua mô hình ngôn ngữ, không phải để huấn luyện một bộ phân lớp mới.
Các khái niệm về tác tử và hệ thống trí tuệ nhân tạo được tham chiếu từ giáo
trình nền tảng~\cite{RussellNorvig2021}.

\subsubsection{Sinh câu hỏi có điều kiện theo nguồn}
Sinh câu hỏi từ tài liệu là quá trình chuyển nội dung văn bản thành câu hỏi
và đáp án dựa trên các ý trong nguồn. Đầu vào của AI ExamHub DAU là nội dung
đã trích xuất cùng môn học, số lượng, loại câu và phân bố độ khó/mức độ tư
duy. Kết quả được chuẩn hoá thành các trường như nội dung câu hỏi, lựa chọn,
đáp án, giải thích và đoạn bằng chứng. Gắn bằng chứng nguồn giúp giảng viên
đối chiếu đề xuất với học liệu trước khi lưu vào ngân hàng.

Khi có cấu hình khóa OpenAI, backend gửi yêu cầu đến dịch vụ mô hình ngôn
ngữ. Đây là suy luận qua API, không phải mô hình do đồ án huấn luyện. Khi
không có khóa hoặc phản hồi không dùng được, backend chuyển sang bộ sinh dự
phòng dựa trên các câu/đoạn trong tài liệu. Bộ dự phòng giúp luồng có thể
hoạt động mà không phụ thuộc dịch vụ ngoài, nhưng không đảm bảo chất lượng
ngữ nghĩa tương đương người biên soạn.

\subsubsection{Mô hình ngôn ngữ và sinh theo chỉ dẫn}
Mô hình ngôn ngữ ước lượng chuỗi từ tiếp theo dựa trên ngữ cảnh đã nhận.
Nhờ huấn luyện trên lượng văn bản lớn, một mô hình có thể làm theo chỉ dẫn
và tạo đầu ra với nhiều dạng cấu trúc; nghiên cứu về mô hình ngôn ngữ few-shot
mô tả khả năng điều kiện hoá đầu ra từ ví dụ trong ngữ cảnh
\cite{Brown2020}. Khả năng tạo câu tự nhiên không tự chứng minh nội dung là
đúng. Mô hình có thể diễn đạt trôi chảy nhưng suy luận sai, bỏ qua điều kiện
hoặc tự bổ sung chi tiết không có căn cứ.

Vì vậy, yêu cầu gửi đến mô hình cần làm rõ vai trò, nội dung nguồn, ma trận
câu hỏi, định dạng đầu ra, quy tắc đáp án và yêu cầu trích nguyên văn đoạn
bằng chứng. Ràng buộc cấu trúc giúp backend có thể phân tích JSON và từ
chối đầu ra sai định dạng, nhưng một JSON hợp lệ vẫn có thể chứa kiến thức
sai. Quy trình sản xuất cần có bước kiểm duyệt chuyên môn trước khi phát
hành.

\subsubsection{Sinh có truy xuất và khả năng truy nguyên}
Một hướng nghiên cứu liên quan là kết hợp mô hình sinh với truy xuất tri
thức bên ngoài. Retrieval-Augmented Generation (RAG) kết nối bộ sinh với
kho tri thức để cung cấp thông tin và dấu vết nguồn cho câu trả lời
\cite{Lewis2020RAG}. Trong nguyên mẫu hiện tại, backend đưa nội dung tài
liệu được người dùng gửi trực tiếp vào yêu cầu, rồi yêu cầu trả
\code{sourceEvidence}; dự án chưa có kho embedding/vector, bộ truy hồi theo
ngữ nghĩa hay pipeline RAG độc lập. Do đó, thuật ngữ “có căn cứ theo tài
liệu” ở đây mô tả điều kiện prompt và trích dẫn nguồn, không khẳng định
kiến trúc RAG đã được triển khai.

\subsubsection{Độ khó và mức độ nhận thức}
Độ khó của câu hỏi không chỉ phụ thuộc độ dài câu. Câu hỏi nhận biết/hiểu
thường yêu cầu giải thích hoặc nhận diện khái niệm; câu vận dụng yêu cầu
dùng quy tắc trong tình huống. Hệ thống lưu mức độ nhận thức
\code{understanding} hoặc \code{application} và các nhãn độ khó, đồng thời
cho phép giảng viên đặt số lượng mong muốn. Các nhãn trong metadata là đầu
vào của quá trình chọn câu hỏi, không phải kết quả đã được hiệu chuẩn theo
một thang đo tâm lý lượng.

\subsubsection{Chất lượng câu hỏi trắc nghiệm}
Câu trắc nghiệm cần có một đáp án đúng duy nhất, phương án nhiễu cùng phạm
trù và không vô lý, lời dẫn rõ ràng, tránh manh mối hình thức và tránh phụ
thuộc vào kiến thức ngoài phạm vi môn học. Bộ sinh của dự án kiểm tra một
số ràng buộc như số lượng/tính khác nhau của lựa chọn và khả năng đối chiếu
bằng chứng. Những kiểm tra này không xác định đầy đủ tính đúng đắn chuyên
môn hay chất lượng sư phạm, nên việc duyệt của giảng viên vẫn là một thành
phần bắt buộc của quy trình sử dụng an toàn.

\subsubsection{Trích xuất nội dung và OCR}
PDF được đọc bằng \code{pdf-parse}; DOC/DOCX được chuyển qua Mammoth, vốn
được thiết kế chủ yếu cho DOCX; TXT được đọc theo UTF-8. Backend cũng có
nhánh OCR ảnh dùng Tesseract.js, khởi tạo với dữ liệu ngôn ngữ tiếng Anh.
Giao diện hiện chỉ cho tải PDF, DOC/DOCX và TXT nên nhánh ảnh chưa phải
luồng tải lên thông thường. Văn bản quét, DOC không tương thích, bố cục phức
tạp hoặc tài liệu tiếng Việt đi qua bộ nhận dạng tiếng Anh có thể làm sai
ký tự và ranh giới câu; nội dung trích xuất cần được rà soát trước khi dùng.

\subsubsection{Mammoth và Tesseract.js}
Mammoth tập trung vào việc chuyển DOCX sang HTML/văn bản, còn Tesseract.js
cung cấp giao diện OCR JavaScript cho ảnh. Việc có thư viện trong backend
không đồng nghĩa chất lượng trích xuất tiếng Việt đã được bảo đảm; chất
lượng cần được đo riêng trên tập tài liệu phù hợp
\cite{MammothDocs,TesseractDocs}.

\subsubsection{Kiểm tra kết quả sinh}
Kết quả mô hình không được đưa thẳng vào đề. Backend chuẩn hoá và kiểm tra
các trường, loại câu, số lựa chọn, tính duy nhất của lựa chọn và đáp án,
mức độ tư duy, độ khó, đồng thời đối chiếu bằng chứng với văn bản nguồn. Nếu
kết quả không đạt điều kiện, hệ thống dùng kết quả dự phòng hoặc loại câu
không hợp lệ. Đây là kiểm tra cấu trúc và một số ràng buộc nội dung; nó
không thay thế việc giảng viên thẩm định kiến thức.

\subsubsection{Mô hình rủi ro và trách nhiệm sử dụng}
Khung quản trị rủi ro AI của NIST tổ chức hoạt động quanh các chức năng
quản trị, lập bản đồ, đo lường và quản lý rủi ro
\cite{NISTAIRMF}. Hướng dẫn UNESCO về AI tạo sinh trong giáo dục nhấn mạnh
cách tiếp cận lấy con người làm trung tâm, độ tuổi phù hợp và bảo vệ dữ
liệu cá nhân \cite{UNESCOGenAI}. Từ góc nhìn hệ thống, rủi ro chính gồm câu
hỏi sai đáp án, dữ liệu nguồn nhạy cảm bị gửi ra ngoài, OCR sai, lạm dụng
tài khoản và phân quyền sai. Biện pháp đề xuất là giảm thiểu dữ liệu gửi
đi, giới hạn quyền, không dùng đầu ra AI làm kết quả cuối cùng nếu chưa
được rà soát, và ghi rõ dịch vụ bên ngoài có thể nhận nội dung tài liệu.

\begin{figure}[H]
\centering
\resizebox{0.96\textwidth}{!}{%
\begin{tikzpicture}[node distance=8mm and 10mm]
  \node[objbox] (source) {Tài liệu\\ PDF/DOCX/TXT};
  \node[clsbox, right=of source] (extract) {Trích xuất\\ văn bản};
  \node[clsbox, right=of extract] (prompt) {Đặt yêu cầu\\ + ma trận đề};
  \node[clsbox, right=of prompt] (generate) {LLM tùy chọn\\ / sinh dự phòng};
  \node[clsbox, below=of generate] (validate) {Kiểm tra cấu trúc,\\ đáp án, bằng chứng};
  \node[objbox, left=of validate] (review) {Giảng viên\\ rà soát / sửa};
  \node[objbox, left=of review] (bank) {Ngân hàng\\ câu hỏi};
  \draw[flowarr] (source) -- (extract);
  \draw[flowarr] (extract) -- (prompt);
  \draw[flowarr] (prompt) -- (generate);
  \draw[flowarr] (generate) -- (validate);
  \draw[flowarr] (validate) -- (review);
  \draw[flowarr] (review) -- (bank);
\end{tikzpicture}}
\caption{Quy trình khai thác tài liệu và rà soát câu hỏi trong hệ thống.}
\label{fig:ai-pipeline}
\end{figure}
Hình~\ref{fig:ai-pipeline} mô tả luồng logic; kết quả của bước kiểm tra
tự động chỉ là cổng cấu trúc ban đầu. Phê duyệt nội dung vẫn thuộc trách
nhiệm giảng viên.

\subsection{Công cụ và thư viện}

\begin{table}[H]
\centering
\caption{Công nghệ được khai báo trong các manifest và cấu hình triển khai.}
\label{tab:congcu}
\begin{tabular*}{\textwidth}{@{\extracolsep{\fill}}>{\bfseries\color{daunavy}}p{0.20\textwidth}>{\color{daumaroon}}p{0.23\textwidth}p{0.47\textwidth}}
\toprule
\rowcolor{gray100}\bfseries\color{daunavy}Thành phần & \bfseries\color{daumaroon}Phiên bản & \bfseries\color{daunavy}Vai trò \\
\midrule
React & 18.3.1 & Xây dựng giao diện web \\
TypeScript & 5.6.3 & Kiểm tra kiểu cho mã giao diện \\
Vite & 5.4.11 & Phát triển và đóng gói frontend \\
Node.js & 20 (Docker) & Môi trường chạy backend \\
Express & 5.1.0 & Xây dựng REST API \\
PostgreSQL & 16 (Docker) & Lưu trữ dữ liệu quan hệ \\
Tesseract.js & 7.x & Thành phần OCR được khai báo trong backend \\
pdf-parse, Mammoth & Theo package.json & Đọc PDF và DOCX \\
Docker Compose & Theo môi trường cài đặt & Khởi chạy các container phát triển \\
\bottomrule
\end{tabular*}
\end{table}

\subsection{Tài liệu nền tảng và công nghệ liên quan}

Các mô hình ngôn ngữ hiện đại có thể sinh văn bản theo chỉ dẫn, nhưng đầu ra
có thể chứa thông tin sai hoặc không bám sát nguồn. Vì vậy, hệ thống ứng
dụng nguyên tắc có điều kiện theo tài liệu, yêu cầu bằng chứng và kiểm tra
đầu ra trước khi giảng viên duyệt; tài liệu API mô hình là nguồn tham khảo
cho cách tích hợp dịch vụ suy luận~\cite{OpenAIAPI}. Nền tảng còn sử dụng
các thành phần ứng dụng web và cơ sở dữ liệu đã được chuẩn hoá trong tài
liệu chính thức của React, Express và PostgreSQL~\cite{ReactDocs,TypeScriptDocs,ViteDocs,ExpressDocs,PostgreSQLDocs}.
Khác với một nghiên cứu mô hình học máy, đồ án tập trung vào tích hợp các
luồng nghiệp vụ và không đưa ra so sánh benchmark giữa các mô hình.

\subsubsection{Cơ sở dữ liệu quan hệ và tính toàn vẹn}
PostgreSQL hỗ trợ khoá chính, khoá ngoại, ràng buộc và giao dịch, phù hợp
với các quan hệ giữa tài khoản, đề, câu hỏi, lượt thi và bài nộp
\cite{PostgreSQLDocs}. Câu hỏi và câu trả lời được biểu diễn dưới dạng JSONB
ở một số cột để giữ lựa chọn và đáp án theo cấu trúc linh hoạt; những liên
kết cần truy vấn và toàn vẹn tham chiếu vẫn có bảng liên kết riêng. Cách
phân chia này tạo thuận lợi cho metadata, nhưng cần tránh để hai biểu diễn
\code{question\_ids} và \code{exam\_questions} sai khác nhau.

\subsubsection{Giao diện thành phần và hợp đồng HTTP}
React tổ chức giao diện thành các component có trạng thái cục bộ và props
\cite{ReactDocs}; TypeScript bổ sung kiểu dữ liệu tĩnh cho mô hình người
dùng, câu hỏi, đề và lượt thi \cite{TypeScriptDocs}. Vite cung cấp máy chủ
phát triển và quy trình build tài nguyên frontend
\cite{ViteDocs}. Backend dùng Node.js và Express để khai báo middleware,
xác thực và tuyến HTTP \cite{NodeDocs,ExpressDocs}. Những công nghệ này
giảm công sức tích hợp, nhưng vẫn cần hợp đồng API nhất quán, xử lý lỗi và
quyền truy cập ở backend.

\subsubsection{Kiểm thử phần mềm và khả năng tái lập}
Bộ kiểm thử backend sử dụng module \code{node:test} có sẵn trong Node.js;
test runner cung cấp cách tổ chức và thực thi các test đồng bộ hoặc bất
đồng bộ \cite{NodeTestDocs}. Compose mô tả dịch vụ backend, frontend và cơ
sở dữ liệu trong một cấu hình để có thể dựng lại stack
\cite{DockerComposeDocs}. Kết quả kiểm thử chỉ có ý nghĩa đối với phiên bản
mã, cấu hình và môi trường đã chạy; vì vậy khi tái lập cần ghi lại commit,
phiên bản Node.js, biến môi trường đã che bí mật và tình trạng cơ sở dữ liệu.

% ---------------------------------------------------------------------
%  CHƯƠNG 2 — PHÂN TÍCH, THIẾT KẾ VÀ TRIỂN KHAI
% ---------------------------------------------------------------------
\chapter{Phân tích, thiết kế và triển khai}
\label{ch:thietke}

\begin{muctieu}
Chương này phân tích vai trò và yêu cầu, mô tả kiến trúc, dữ liệu, quy
trình sinh câu hỏi và tổ chức thi, rồi trình bày cài đặt và thử nghiệm.
\end{muctieu}

\section{Phân tích bài toán và yêu cầu}

Hệ thống có ba vai trò chính: quản trị viên quản lý tài khoản và dữ liệu
toàn hệ thống; giảng viên quản lý câu hỏi, tài liệu, đề thi, lớp học và
chấm bài; sinh viên tham gia lớp/phòng thi và xem lịch sử dự thi. Các yêu
cầu chức năng trọng tâm được tóm tắt trong Bảng~\ref{tab:yeucau}.

\begin{table}[H]
\centering
\caption{Nhóm yêu cầu chức năng của hệ thống.}
\label{tab:yeucau}
\begin{tabular*}{\textwidth}{@{\extracolsep{\fill}}p{0.12\textwidth}p{0.31\textwidth}p{0.49\textwidth}}
\toprule
\rowcolor{gray100}\bfseries\color{daunavy}Mã & \bfseries\color{daunavy}Chức năng & \bfseries\color{daunavy}Mô tả \\
\midrule
FR1 & Tài khoản và phân quyền & Đăng nhập, hồ sơ người dùng và giới hạn chức năng theo vai trò. \\
FR2 & Ngân hàng câu hỏi & Tạo, sửa, xoá, lọc và phân loại câu hỏi theo môn, loại và độ khó. \\
FR3 & Hỗ trợ tạo đề & Đọc tài liệu, đề xuất câu hỏi, chọn phân bố, lưu đề và tạo mã đề. \\
FR4 & Duyệt đề & Chuyển đề qua các trạng thái nháp, chờ duyệt, đã duyệt và phát hành. \\
FR5 & Thi trực tuyến & Tạo phòng/mã tham gia, lưu câu trả lời, tiếp tục lượt thi và tự chấm trắc nghiệm. \\
FR6 & Lớp học & Tạo lớp, tham gia bằng mã, quản lý trạng thái buổi học và thành viên. \\
FR7 & Chấm và thống kê & Giảng viên chấm bài, ghi nhận phản hồi và xem số liệu theo quyền. \\
\bottomrule
\end{tabular*}
\end{table}

Các yêu cầu phi chức năng gồm bảo vệ API bằng xác thực và phân quyền, lưu
dữ liệu bền vững qua PostgreSQL, kiểm tra đầu vào ở các luồng quan trọng,
và hỗ trợ chạy các thành phần chính bằng Docker Compose. Báo cáo chưa xác
nhận ngưỡng thời gian phản hồi hay số người dùng đồng thời; các tiêu chí đó
cần được đo riêng trước triển khai.

\subsection{Tác nhân và kịch bản sử dụng}
Các tác nhân tương tác với hệ thống qua giao diện web. Quản trị viên tạo và
quản lý tài khoản, phân quyền và xem số liệu toàn hệ thống. Giảng viên nạp
học liệu, xem các câu hỏi được đề xuất, hiệu chỉnh câu hỏi, lưu thành ngân
hàng, tạo/duyệt đề, mở phòng thi, giám sát, chấm bài và xem phân tích.
Sinh viên đăng nhập hoặc tham gia lớp/phòng thi theo quyền, trả lời câu
hỏi, đánh dấu nội dung cần xem lại và nộp bài. Một lượt dự thi được gắn với
phòng, đề và tài khoản sinh viên.

Bảng~\ref{tab:usecase} mô tả một số kịch bản trọng yếu. Kịch bản thể hiện
hành vi mong muốn của sản phẩm; các bước được triển khai đến đâu được đánh
giá riêng trong Chương~\ref{ch:ketqua}.

\begin{table}[H]
\centering
\caption{Đặc tả các kịch bản sử dụng chính.}
\label{tab:usecase}
\begin{tabularx}{\textwidth}{>{\bfseries\color{daunavy}}p{0.16\textwidth} X X}
\toprule
\rowcolor{gray100}Kịch bản & Luồng chính & Kết quả / ngoại lệ cần xử lý \\
\midrule
UC1 -- Sinh câu hỏi & Giảng viên chọn tài liệu, môn học, số lượng, dạng câu và phân bố; gửi phân tích; xem ứng viên, bằng chứng và hiệu chỉnh. & Ứng viên có nguồn để đối chiếu; từ chối tài liệu rỗng, lỗi đọc hoặc đầu ra không hợp lệ. \\
UC2 -- Tạo/phát hành đề & Chọn câu hỏi; khai báo thời lượng, thang điểm và metadata; lưu nháp; gửi duyệt; phê duyệt và phát hành. & Đề lưu kèm liên kết câu hỏi; lỗi ma trận không đủ câu phải được thông báo. \\
UC3 -- Tổ chức thi & Giảng viên chọn đề, đặt thời gian, tạo mã; sinh viên nhập mã, làm bài, lưu và nộp. & Một sinh viên tiếp tục đúng lượt thi đã bắt đầu; bài nộp được tự chấm trắc nghiệm. \\
UC4 -- Chấm tự luận & Giảng viên chọn bài đã nộp, nhập điểm trong giới hạn và phản hồi, lưu kết quả. & Kết quả gắn với bài nộp; không cho điểm vượt tổng điểm đề. \\
UC5 -- Tham gia lớp & Giảng viên tạo lớp, chọn kiểu kết nối; sinh viên nhập mã; giảng viên mở/đóng buổi học. & Thành viên được lưu; lớp đã kết thúc hoặc đóng nhận thành viên thì từ chối lượt vào mới. \\
\bottomrule
\end{tabularx}
\end{table}

\begin{figure}[H]
\centering
\resizebox{0.9\textwidth}{!}{%
\begin{tikzpicture}[node distance=8mm and 12mm]
  \node[objbox] (teacher) {Giảng viên};
  \node[objbox, below=of teacher] (student) {Sinh viên};
  \node[objbox, above=of teacher] (admin) {Quản trị viên};
  \node[clsbox, right=of teacher] (bank) {Quản lý ngân hàng\\câu hỏi};
  \node[clsbox, right=of bank] (exam) {Tạo / duyệt\\đề thi};
  \node[clsbox, right=of exam] (room) {Thi và giám sát\\phòng thi};
  \node[clsbox, below=of exam] (classroom) {Lớp học\\trực tuyến};
  \node[clsbox, below=of bank] (grading) {Chấm bài\\và thống kê};
  \draw[flowarr] (teacher) -- (bank);
  \draw[flowarr] (teacher) -- (exam);
  \draw[flowarr] (teacher) -- (room);
  \draw[flowarr] (teacher) -- (grading);
  \draw[flowarr] (teacher) -- (classroom);
  \draw[flowarr] (student) -- (room);
  \draw[flowarr] (student) -- (classroom);
  \draw[flowarr] (student) -- (grading);
  \draw[flowarr] (admin) -- (bank);
  \draw[flowarr] (admin) -- (exam);
  \draw[flowarr] (admin) -- (room);
\end{tikzpicture}}
\caption{Quan hệ giữa các vai trò và nhóm chức năng ở mức use case.}
\label{fig:use-case}
\end{figure}
Hình~\ref{fig:use-case} chỉ thể hiện nhóm chức năng theo vai trò, không thay
thế bảng endpoint hoặc chính sách phân quyền chi tiết.

\subsection{Yêu cầu phi chức năng và tiêu chí nghiệm thu}
Yêu cầu phi chức năng được phân thành an toàn, toàn vẹn dữ liệu, khả dụng,
dễ vận hành và khả năng tiếp cận. Trạng thái đạt/đo được trình bày ở
Bảng~\ref{tab:nfr}; các ngưỡng chưa đo không được xem là đã thoả mãn.
\begin{table}[H]
\centering
\caption{Yêu cầu phi chức năng và mức xác minh hiện có.}
\label{tab:nfr}
\begin{tabularx}{\textwidth}{>{\bfseries\color{daunavy}}p{0.20\textwidth} X p{0.23\textwidth}}
\toprule
\rowcolor{gray100}Thuộc tính & Tiêu chí yêu cầu & Bằng chứng hiện có \\
\midrule
Phân quyền & Mỗi thao tác ghi dữ liệu chỉ được phép với vai trò và chủ sở hữu thích hợp. & Có middleware và test phân quyền; chưa có kiểm thử độc lập đầy đủ. \\
Toàn vẹn & Khóa ngoại, ràng buộc, giao dịch tạo đề/phòng và tính duy nhất. & Có schema, migration và transaction ở các luồng chính. \\
Khả dụng & Lỗi dịch vụ ngoài không được biến thành kết quả thành công giả. & Có fallback cho sinh nội dung; cần kiểm tra log và UX ở triển khai thật. \\
Hiệu năng & Giao diện và API đáp ứng thời gian chấp nhận được ở tải mục tiêu. & Chưa có ngưỡng hoặc đo tải. \\
Khả năng tiếp cận & Điều khiển bằng bàn phím, nhãn biểu mẫu, trạng thái và tương phản phù hợp. & Chưa thực hiện đánh giá WCAG có hệ thống~\cite{WCAG22}. \\
Quản lý dữ liệu & Hạn chế dữ liệu cá nhân/tài liệu gửi cho bên thứ ba; quy định lưu trữ và xoá. & Chưa có kiểm toán quyền riêng tư hoặc chính sách vận hành hoàn chỉnh. \\
\bottomrule
\end{tabularx}
\end{table}

\section{Thiết kế kiến trúc hệ thống}

Kiến trúc triển khai được trình bày ở Hình~\ref{fig:kientruc}. Frontend
React gửi yêu cầu HTTP đến REST API Express; backend áp dụng xác thực,
nghiệp vụ và truy vấn PostgreSQL. Luồng tài liệu đi qua bộ trích xuất văn
bản, sau đó có thể gọi dịch vụ mô hình ngôn ngữ ngoài hoặc bộ sinh dự phòng.
Các thư mục \code{ai-services/} và \code{ml-models/} hiện là khung tổ chức,
không phải các dịch vụ độc lập đang được gọi trong luồng runtime.

\begin{figure}[H]
\centering
\resizebox{\textwidth}{!}{%
\begin{tikzpicture}[node distance=12mm and 14mm]
  \node[objbox] (user) {Giảng viên /\\ sinh viên / admin};
  \node[clsbox, right=of user] (web) {Frontend\\ React + TypeScript};
  \node[clsbox, right=of web] (api) {REST API\\ Node.js + Express};
  \node[clsbox, right=of api] (db) {PostgreSQL\\ users, exams, attempts};
  \node[clsbox, below=of api] (doc) {Trích xuất văn bản\\ PDF / DOCX / OCR};
  \node[objbox, right=of doc] (ai) {OpenAI API tùy chọn\\ hoặc bộ sinh dự phòng};
  \draw[flowarr] (user) -- (web);
  \draw[flowarr] (web) -- (api);
  \draw[flowarr] (api) -- (db);
  \draw[flowarr] (api) -- (doc);
  \draw[flowarr] (doc) -- (ai);
  \draw[flowarr] (ai.north) |- (api.south);
\end{tikzpicture}}
\caption{Kiến trúc logic và luồng xử lý chính của AI ExamHub DAU.}
\label{fig:kientruc}
\end{figure}

\subsection{Phân rã thành phần}
Frontend đóng vai trò client hiển thị view theo vai trò và gọi các hàm
trong lớp dịch vụ API. Backend tập trung xử lý các quy tắc miền, xác thực
JWT, kiểm tra vai trò, thao tác cơ sở dữ liệu và tích hợp mô hình ngoài.
PostgreSQL giữ dữ liệu nghiệp vụ; Docker Compose nối backend, frontend và
cơ sở dữ liệu trong môi trường phát triển. Các module OCR/RAG/embedding
trong cây thư mục dự án là vị trí quy hoạch cho tương lai, không đồng nghĩa
các container AI này đã tồn tại trong cấu hình Compose hiện tại.

\subsection{Luồng xử lý phân tích tài liệu}
Luồng phân tích nhận \code{documentId}, môn học, số lượng, nội dung, ma
trận loại câu, phân bố độ khó và mức độ nhận thức. Backend từ chối nội
dung rỗng; nếu có nội dung, bộ phân tích tách các câu hỏi được định dạng
sẵn, đồng thời gọi bộ sinh từ tài liệu. Các ứng viên được hợp nhất, chuẩn
hoá và lọc theo ma trận. Những câu phù hợp được lưu vào ngân hàng câu hỏi
trước khi trả về giao diện. Mô hình dữ liệu do đó phân biệt việc phân tích
đề xuất với việc người dùng chọn câu để tạo đề hoàn chỉnh.

\subsection{Luồng thi trực tuyến}
Giảng viên tạo phòng từ một hoặc nhiều đề đã lưu. Mã phòng được tạo ngẫu
nhiên; sinh viên dùng mã để bắt đầu lượt thi. Backend ghi một lượt cho
cặp phòng--sinh viên, bộ câu hỏi được giao, đáp án, câu đánh dấu, vị trí
hiện tại, thời hạn và thời điểm nộp. Khi nộp, bài thi trắc nghiệm được
tự chấm; đề tự luận có thể yêu cầu chấm riêng. Giảng viên có thể xem tiến
độ và sự kiện chuyển tab được ghi lại. Sự kiện chuyển tab là tín hiệu hỗ
trợ giám sát, không đủ để tự kết luận có gian lận.

\section{Thiết kế dữ liệu}

Hệ thống không sử dụng tập dữ liệu huấn luyện phân lớp. Dữ liệu nghiệp vụ
gồm tài khoản, câu hỏi, đề thi, lượt dự thi, bài nộp, kết quả và lớp học.
Các bảng lõi được mô tả khái quát ở Bảng~\ref{tab:dulieu}; các trường chi
tiết được khai báo trong lược đồ PostgreSQL của dự án.

\begin{table}[H]
\centering
\caption{Các thực thể chính trong cơ sở dữ liệu.}
\label{tab:dulieu}
\begin{tabular*}{\textwidth}{@{\extracolsep{\fill}}>{\bfseries\color{daunavy}}p{0.19\textwidth}>{\color{daumaroon}}p{0.28\textwidth}p{0.43\textwidth}}
\toprule
\rowcolor{gray100}\bfseries\color{daunavy}Thực thể & \bfseries\color{daumaroon}Bảng & \bfseries\color{daunavy}Vai trò \\
\midrule
Tài khoản & users & Email, vai trò, hồ sơ và thông tin định danh. \\
Câu hỏi & questions & Nội dung, lựa chọn, đáp án, bằng chứng nguồn và rubric. \\
Đề thi & exams, exam\_questions & Metadata đề và liên kết có thứ tự tới câu hỏi. \\
Thi trực tuyến & exam\_rooms, exam\_attempts & Phòng thi, câu hỏi được giao, đáp án và thời điểm. \\
Bài nộp/kết quả & submissions, results & Trạng thái nộp, điểm số và phản hồi. \\
Lớp học & classroom\_rooms, classroom\_members & Buổi học, mã tham gia và thành viên. \\
\bottomrule
\end{tabular*}
\end{table}

\subsection{Sơ đồ thực thể--liên kết}
Các quan hệ lõi trong schema được tóm tắt tại Hình~\ref{fig:erd}. Người
dùng tạo câu hỏi/đề/phòng/lớp; đề liên kết đến câu hỏi qua bảng nối; lượt
thi thuộc về một sinh viên, một phòng và một đề; bài nộp liên kết với kết
quả. Bảng lớp học nối người dùng sinh viên với phòng học bằng khoá ghép.
\begin{figure}[H]
\centering
\resizebox{\textwidth}{!}{%
\begin{tikzpicture}[node distance=10mm and 13mm]
  \node[clsbox] (users) {users};
  \node[clsbox, right=of users] (questions) {questions};
  \node[clsbox, right=of questions] (exams) {exams};
  \node[objbox, below=of questions] (examq) {exam\_questions};
  \node[clsbox, below=of exams] (rooms) {exam\_rooms};
  \node[objbox, right=of rooms] (roomexams) {exam\_room\_exams};
  \node[clsbox, below=of rooms] (attempts) {exam\_attempts};
  \node[clsbox, left=of attempts] (submissions) {submissions};
  \node[clsbox, left=of submissions] (results) {results};
  \node[objbox, below=of users] (rubrics) {rubrics};
  \node[clsbox, below=of attempts] (classrooms) {classroom\_rooms};
  \node[objbox, left=of classrooms] (members) {classroom\_members};
  \draw[flowarr] (users) -- (questions);
  \draw[flowarr] (users) |- (exams);
  \draw[flowarr] (questions) -- (examq);
  \draw[flowarr] (exams) -- (examq);
  \draw[flowarr] (exams) -- (rooms);
  \draw[flowarr] (rooms) -- (roomexams);
  \draw[flowarr] (exams) -- (roomexams);
  \draw[flowarr] (rooms) -- (attempts);
  \draw[flowarr] (exams) -- (attempts);
  \draw[flowarr] (users) |- (attempts);
  \draw[flowarr] (attempts) -- (submissions);
  \draw[flowarr] (submissions) -- (results);
  \draw[flowarr] (exams) |- (rubrics);
  \draw[flowarr] (users) |- (classrooms);
  \draw[flowarr] (classrooms) -- (members);
  \draw[flowarr] (users) |- (members);
\end{tikzpicture}}
\caption{Sơ đồ quan hệ giữa các nhóm bảng nghiệp vụ chính.}
\label{fig:erd}
\end{figure}

\section{Thiết kế mô hình hoặc cơ sở tri thức}

Thành phần tạo câu hỏi có hai nhánh. Nhánh mô hình ngôn ngữ được gọi qua
OpenAI API khi cấu hình \code{OPENAI\_API\_KEY}; yêu cầu đầu ra gồm loại câu,
nội dung, lựa chọn, đáp án, giải thích và bằng chứng nguồn. Backend kiểm tra
định dạng, ràng buộc lựa chọn và mức độ bám nguồn trước khi nhận câu. Nhánh
dự phòng tách nội dung thành các đoạn/câu, nhận diện các mẫu kiến thức và
tạo câu hỏi theo quy tắc; nhánh này giữ hệ thống hoạt động khi dịch vụ ngoài
không sẵn sàng nhưng có phạm vi suy luận hạn chế.

Ở giai đoạn tạo đề, câu hỏi được lọc theo loại, độ khó và mức độ tư duy rồi
lưu cùng metadata đề trong cơ sở dữ liệu. Khi tạo nhiều mã đề, backend xáo
trộn thứ tự câu/lựa chọn và ánh xạ lại đáp án. Trong lúc thi, mỗi sinh viên
được gắn với một bộ câu hỏi cho lượt thi, câu trả lời và vị trí hiện tại
được lưu để hỗ trợ tiếp tục bài làm. Trắc nghiệm được tự chấm theo đáp án;
tự luận được đưa qua quy trình giảng viên chấm và ghi phản hồi.

\subsection{Thiết kế phân quyền}
Chính sách quyền được áp dụng theo vai trò và một số trường hợp theo quyền
sở hữu bản ghi. Admin quản lý dữ liệu toàn hệ thống; giảng viên thực hiện
chức năng giáo viên và tài nguyên do mình quản lý; sinh viên chỉ thao tác
với lượt thi/lớp mà mình được phép truy cập. Một token xác thực chỉ xác
nhận danh tính; backend còn phải kiểm tra quyền của tài nguyên cụ thể. RFC
7519 định nghĩa cấu trúc JWT, nhưng việc sử dụng token an toàn phụ thuộc
vào thuật toán, thời hạn, kiểm tra claim, bảo vệ khóa và thu hồi phiên
\cite{RFC7519}. Đăng nhập Google cần xác thực luồng OAuth/OpenID Connect,
kiểm tra trạng thái/redirect và liên kết tài khoản đúng theo danh tính đã
xác minh \cite{GoogleOIDC}.

\subsection{Bảo vệ API và thành phần AI}
OWASP Top 10 và OWASP API Security Top 10 nêu các nhóm rủi ro như kiểm soát
truy cập sai, lỗi xác thực, cấu hình không an toàn và khai thác tài nguyên
\cite{OWASPTop10,OWASPAPI}. Với mô hình ngôn ngữ, OWASP LLM Top 10 nhấn mạnh
prompt injection, rò rỉ thông tin và đầu ra không được kiểm soát
\cite{OWASPLLM}. Ở hệ thống thi, chính sách thực tế cần kiểm tra quyền sở
hữu từng đề/phòng, không trả đáp án đúng cho sinh viên trong khi thi, giới
hạn kích thước/định dạng tải lên, lưu khóa API ở secret store, giảm thiểu
nội dung nhạy cảm gửi ra mô hình và xem đầu ra mô hình như dữ liệu không
tin cậy. Phân tích này là yêu cầu thiết kế và không phải kết quả kiểm thử
xâm nhập.

% ---------------------------------------------------------------------
%  NỘI DUNG CHƯƠNG 2 — CÀI ĐẶT VÀ THỬ NGHIỆM
% ---------------------------------------------------------------------
\section{Cài đặt và thử nghiệm}
\label{sec:caidat}

\begin{muctieu}
Phần này mô tả công nghệ, các thành phần đã triển khai và cách kiểm thử
những luồng nghiệp vụ chính của hệ thống.
\end{muctieu}

\subsection{Môi trường triển khai}

Dự án được tổ chức thành frontend web, backend API, cơ sở dữ liệu và cấu
hình Docker Compose. Frontend dùng React 18 và TypeScript 5; backend dùng
Node.js với Express 5; PostgreSQL 16 là dịch vụ cơ sở dữ liệu trong cấu
hình container. Các phiên bản thư viện chính được liệt kê ở
Bảng~\ref{tab:congcu}. Dockerfile của frontend dùng bước build Vite và Nginx
để phục vụ tệp tĩnh; backend chạy trên ảnh Node.js Alpine.

Mã nguồn được kiểm tra tại thư mục cục bộ
\code{E:\textbackslash AI-EXAMHUB-DAU}. Cấu hình môi trường được cung cấp qua
các tệp \code{.env.example}; khóa dịch vụ và mật khẩu dùng khi triển khai
thực tế cần được thay bằng giá trị riêng, không đưa vào báo cáo hoặc kho mã
nguồn. Hướng dẫn chạy từ mã nguồn được trình bày ở Phụ lục.

\subsubsection{Cấu trúc mã nguồn}
Các thành phần được tách theo trách nhiệm. \code{frontend/web/src/pages/}
chứa các màn hình nghiệp vụ; \code{frontend/web/src/services/api.ts} là lớp
gọi backend. \code{backend/api/src/app.js} khởi tạo ứng dụng Express và các
route; \code{backend/api/src/routes/} chứa các nhóm API. SQL nền tảng nằm
trong \code{database/schema.sql}; các thay đổi theo phiên bản được đặt tại
\code{database/migrations/}.

\begin{table}[H]
\centering
\caption{Phân chia trách nhiệm theo thư mục mã nguồn.}
\label{tab:modules}
\begin{tabularx}{\textwidth}{>{\bfseries\color{daunavy}}p{0.29\textwidth} X}
\toprule
\rowcolor{gray100}Đường dẫn & Trách nhiệm trong hệ thống \\
\midrule
\code{frontend/web/src/pages/} & Màn hình tạo đề, câu hỏi, thi, lớp học, người dùng, chấm và thống kê. \\
\code{frontend/web/src/services/api.ts} & Tạo yêu cầu HTTP, giữ kiểu dữ liệu API và quản lý thông tin xác thực ở phía client. \\
\code{backend/api/src/app.js} & Khởi tạo Express, tích hợp phân tích tài liệu và đăng ký route. \\
\code{backend/api/src/routes/} & Xác thực, câu hỏi, đề thi, phòng thi, lớp học, người dùng, bài nộp và phân tích. \\
\code{backend/api/src/middlewares/} & Xác thực JWT và kiểm tra vai trò. \\
\code{database/schema.sql} & Định nghĩa bảng, khoá, ràng buộc và chỉ mục chính. \\
\code{database/migrations/} & Chuyển đổi schema cho các phiên bản chức năng. \\
\code{docker-compose.yml} & Khai báo dịch vụ phát triển và mạng/kết nối giữa các container. \\
\bottomrule
\end{tabularx}
\end{table}

\subsubsection{Biến cấu hình và phụ thuộc ngoài}
Bảng~\ref{tab:config} mô tả nhóm cấu hình theo mục đích, không nêu giá trị
bí mật. Ứng dụng có thể chạy phần sinh dự phòng mà không cấu hình API mô
hình; chức năng Google OAuth cũng phụ thuộc thông tin client do nhà triển
khai cấp. Hướng dẫn Node.js và Express được dùng để tham chiếu hành vi môi
trường chạy và định tuyến \cite{NodeDocs,ExpressDocs}.
\begin{table}[H]
\centering
\caption{Nhóm cấu hình cần thiết khi chạy hệ thống.}
\label{tab:config}
\begin{tabularx}{\textwidth}{>{\bfseries\color{daunavy}}p{0.22\textwidth}p{0.33\textwidth}X}
\toprule
\rowcolor{gray100}Nhóm biến & Ví dụ tên biến & Mục đích / lưu ý \\
\midrule
API và xác thực & \code{PORT}, \code{JWT\_SECRET}, \code{JWT\_EXPIRES\_IN} & Cổng API, ký và thời hạn token; secret phải duy nhất, không dùng giá trị dev khi triển khai. \\
Cơ sở dữ liệu & \code{DB\_HOST}, \code{DB\_PORT}, \code{DB\_NAME}, \code{DB\_USER}, \code{DB\_PASSWORD} & Kết nối PostgreSQL; mật khẩu lấy từ môi trường an toàn. \\
Giao diện & \code{VITE\_API\_URL} & Địa chỉ gốc của API được đóng vào frontend. \\
LLM tùy chọn & \code{OPENAI\_API\_KEY}, \code{OPENAI\_MODEL} & Gọi API sinh nội dung; khóa không được chuyển xuống trình duyệt. \\
Đăng nhập Google & \code{GOOGLE\_CLIENT\_ID}, \code{GOOGLE\_CLIENT\_SECRET}, \code{GOOGLE\_REDIRECT\_URI} & Chỉ hoạt động khi được đăng ký/cấu hình đúng tại nhà cung cấp. \\
Nguồn gọi API & \code{CORS\_ORIGIN}, \code{FRONTEND\_URL} & Chính sách CORS và URL chuyển hướng OAuth. \\
\bottomrule
\end{tabularx}
\end{table}

\subsubsection{Sơ đồ triển khai các thành phần}
Hình~\ref{fig:deployment} khái quát luồng kết nối của phiên bản hiện tại.
Trình duyệt tải giao diện tĩnh do Nginx phục vụ, sau đó gửi yêu cầu đến API
Express; backend truy cập PostgreSQL và chỉ gọi nhà cung cấp LLM bên ngoài
khi đã cấu hình khóa. Các thành phần frontend, backend và cơ sở dữ liệu được
khai báo trong Docker Compose \cite{DockerComposeDocs}; mô hình không hàm ý
dịch vụ AI riêng biệt chạy trong Compose.
\begin{figure}[H]
\centering
\resizebox{0.95\textwidth}{!}{%
\begin{tikzpicture}[node distance=10mm and 13mm]
  \node[objbox] (browser) {Trình duyệt\\React};
  \node[clsbox, right=of browser] (web) {Nginx\\tệp giao diện};
  \node[clsbox, right=of web] (api) {Express API\\nghiệp vụ};
  \node[clsbox, right=of api] (db) {PostgreSQL\\dữ liệu};
  \node[clsbox, below=of api] (llm) {Dịch vụ LLM\\tuỳ cấu hình};
  \draw[flowarr] (browser) -- node[above,font=\small] {HTTPS} (web);
  \draw[flowarr] (web) -- node[above,font=\small] {HTTP API} (api);
  \draw[flowarr] (api) -- node[above,font=\small] {SQL} (db);
  \draw[flowarr,dashed] (api) -- node[right,font=\small] {HTTPS} (llm);
\end{tikzpicture}}
\caption{Sơ đồ triển khai logic và các kết nối chính.}
\label{fig:deployment}
\end{figure}

\subsection{Các thành phần chính}

\subsubsection{Giao diện web}
Ứng dụng React có màn hình đăng nhập, trang tổng quan, ngân hàng câu hỏi,
trình tạo đề, duyệt đề, quản lý người dùng, chấm bài, phân tích, thi trực
tuyến và lớp học. Giao diện gọi các hàm tập trung trong dịch vụ API ở
\code{frontend/web/src/services/api.ts}. Một số trang điều phối luồng cụ thể
gồm \code{ExamBuilderPage.tsx}, \code{OnlineExamPage.tsx},
\code{QuestionBankPage.tsx} và \code{GradingPage.tsx}.

\subsubsection{API và nghiệp vụ}
Backend Express tổ chức endpoint theo nhóm xác thực, người dùng, câu hỏi,
đề thi, bài nộp, phân tích, phòng thi và lớp học. Middleware xác thực token
và kiểm tra vai trò được dùng để giới hạn thao tác. Endpoint phân tích tài
liệu tiếp nhận nội dung, gọi dịch vụ ngôn ngữ nếu cấu hình, kiểm tra đầu ra
và trả các câu hỏi đề xuất kèm thông tin bằng chứng. Nhóm endpoint thi trực
tuyến hỗ trợ tạo phòng, bắt đầu/tiếp tục lượt thi, lưu câu trả lời, nộp bài
và giám sát tiến độ.

\subsubsection{Lưu trữ và xử lý tài liệu}
Lược đồ PostgreSQL lưu người dùng, câu hỏi, đề thi, phòng/lượt thi, bài nộp,
kết quả và thành viên lớp. Các thay đổi lược đồ được quản lý bằng các tệp
migration trong \code{database/migrations/}. Tệp tải lên được nhận trong bộ
nhớ của backend để trích xuất nội dung; PDF dùng \code{pdf-parse}, DOC/DOCX
được chuyển qua Mammoth và TXT được giải mã UTF-8. Backend có nhánh OCR ảnh
dùng Tesseract.js với dữ liệu tiếng Anh, nhưng giao diện hiện không cung cấp
đầu vào ảnh. Báo cáo không giả định tài liệu đã được lưu trữ lâu dài hay
mọi tài liệu được trích xuất chính xác.

\subsubsection{Quản lý vòng đời đề thi}
Đề lưu trạng thái nghiệp vụ nháp, chờ duyệt, đã duyệt và phát hành cùng
metadata như thời lượng, tổng điểm, năm học, học kỳ và ma trận câu hỏi. Khi
lưu đề, backend kiểm tra số lượng, loại câu, điểm và thời lượng; các thao
tác cần nhiều câu lệnh sử dụng transaction để tránh lưu dở dữ liệu giữa
bảng đề và bảng liên kết. Câu hỏi thuộc đề được truy vấn qua
\code{exam\_questions} theo vị trí, thay vì dựa duy nhất vào một danh sách
ID không có ràng buộc quan hệ.

\subsubsection{Quản lý phiên làm bài}
Khi sinh viên vào phòng, API tạo hoặc lấy lại lượt thi. API lưu đáp án theo
ID câu hỏi, danh sách câu đánh dấu, chỉ số câu hiện tại, thời hạn và thời
điểm nộp. Khi nộp, bài thi trắc nghiệm được tự chấm; đề tự luận có thể yêu
cầu chấm riêng. Trước khi trả câu hỏi cho sinh viên, route loại các trường
đáp án đúng, đáp án tham khảo và bằng chứng nội bộ. Cách tách thông tin
đáp án là kiểm soát quan trọng vì ẩn đáp án chỉ ở giao diện không đủ để
bảo vệ bài thi.

\subsubsection{Thống kê và chấm điểm}
Trang thống kê hiển thị số liệu theo quyền người dùng; quản trị viên có
phạm vi tổng thể, giảng viên tập trung dữ liệu liên quan đề/lớp của mình,
còn sinh viên xem lượt dự thi. Chấm thủ công nhận điểm và phản hồi, sau đó
lưu vào bảng \code{results}; điểm được kiểm tra không âm và không vượt tổng
điểm của đề. Giao diện có nội dung mô tả “AI feedback”, nhưng backend được
kiểm tra trong đồ án không chứng minh đã có pipeline tự chấm tự luận bằng
mô hình; báo cáo chỉ ghi nhận luồng chấm do giảng viên thực hiện.

\subsection{Giao diện và các luồng chính}

Các màn hình triển khai bao phủ hai nhóm luồng. Giảng viên có thể quản lý
câu hỏi, phân tích tài liệu, chọn câu để tạo đề, gửi duyệt, phát hành và
tạo phòng thi. Sinh viên đăng nhập, tham gia bằng mã, làm bài, lưu tiến độ,
tiếp tục lượt thi và xem kết quả. Trang giám sát của giảng viên tải lại
tiến độ theo chu kỳ; trang lớp học quản lý mã tham gia và danh sách thành
viên. Báo cáo hiện không kèm ảnh chụp giao diện vì thư mục dự án không có
ảnh chụp màn hình sản phẩm phù hợp để làm minh chứng.

\subsubsection{Minh hoạ bố cục khu vực tạo đề}
Hình~\ref{fig:wireframe} là sơ đồ bố cục khái niệm được dựng để giải thích
các vùng chức năng, không phải ảnh chụp từ trình duyệt. Nó chỉ ra các bước
người dùng cần quan sát: cấu hình đề, nạp nguồn, chọn câu hỏi và xác nhận
trước khi lưu.
\begin{figure}[H]
\centering
\resizebox{0.92\textwidth}{!}{%
\begin{tikzpicture}[x=1cm,y=1cm]
  \draw[rounded corners=3pt, line width=1pt, color=daunavy] (0,0) rectangle (15,8);
  \fill[lightnavy] (0,7) rectangle (15,8);
  \node[anchor=west,color=daunavy,font=\bfseries] at (0.4,7.5) {AI EXAMHUB DAU \hspace{1em} Tạo đề thi};
  \fill[gray100] (0.4,0.5) rectangle (4.4,6.6);
  \node[anchor=west,font=\bfseries,color=daunavy] at (0.7,6.15) {1. Cấu hình đề};
  \node[anchor=west,align=left] at (2.4,5.25) {Tên đề / môn học\\Thời lượng / tổng điểm\\Loại câu / độ khó\\Mức độ nhận thức};
  \draw[color=gray300] (0.7,3.8) -- (4.1,3.8);
  \node[anchor=west,align=left] at (2.4,3.0) {Chọn tài liệu nguồn\\Tải PDF / DOCX / TXT};
  \node[anchor=west,font=\bfseries,color=daunavy] at (0.7,1.35) {Phân tích tài liệu};
  \draw[rounded corners=2pt,color=daumaroon] (4.8,4.4) rectangle (14.5,6.6);
  \node[anchor=west,font=\bfseries,color=daunavy] at (5.2,6.15) {2. Nội dung nguồn và trích dẫn};
  \node[anchor=west,align=left] at (9.6,5.35) {Xem văn bản đã trích xuất\\Đối chiếu đoạn bằng chứng\\Điều chỉnh thông tin môn học};
  \draw[rounded corners=2pt,color=daumaroon] (4.8,1.0) rectangle (14.5,4.0);
  \node[anchor=west,font=\bfseries,color=daunavy] at (5.2,3.55) {3. Danh sách câu hỏi đề xuất};
  \node[anchor=west,align=left] at (9.6,2.4) {Câu hỏi trắc nghiệm / tự luận\\Mức độ, đáp án và căn cứ\\Chọn, bỏ hoặc hiệu chỉnh câu};
  \fill[lightwarm] (11.4,0.3) rectangle (14.5,0.8);
  \node[anchor=center,font=\bfseries,color=daumaroon] at (12.95,0.55) {Lưu đề};
\end{tikzpicture}}
\caption{Wireframe khái niệm cho các khu vực của quy trình tạo đề.}
\label{fig:wireframe}
\end{figure}

\subsubsection{Trạng thái thông báo và lỗi}
Trong quy trình tạo tài liệu, trạng thái giao diện cần phân biệt đang tải,
trích xuất xong, không có nội dung đọc được, không sinh được ứng viên, đã
sinh được ứng viên và lỗi dịch vụ. Người dùng phải biết liệu câu hỏi thuộc
nhánh LLM hay nhánh trích xuất dự phòng nếu hệ thống trả trường đó. Lỗi
không nên hiển thị dưới dạng thông báo thành công, và thao tác thử lại
không được tạo bản ghi trùng ngoài ý muốn. Những tiêu chí này được đưa vào
ma trận kiểm thử đề xuất ở Phụ lục.

\subsection{Kiểm thử}

Backend có bộ kiểm thử tích hợp tại \code{backend/api/tests/api.test.js}.
Trong lần chạy phục vụ báo cáo, toàn bộ 35 bài kiểm thử đạt, không có bài
thất bại. Các trường hợp tập trung vào đăng nhập và phân quyền, CRUD tài
khoản/câu hỏi/đề thi, phân tích tài liệu, sinh câu hỏi khi không cấu hình
mô hình ngoài, phòng thi trực tuyến, tiếp tục bài làm, tự chấm trắc nghiệm
và quản lý lớp học.

Frontend được kiểm tra bằng lệnh \code{npm run build}; TypeScript hoàn tất
kiểm tra kiểu và Vite tạo gói production thành công. Kết quả này xác nhận
build và các hành vi được kiểm thử, không phải đánh giá độ chính xác của AI,
độ chính xác OCR, tải đồng thời hay thời gian phản hồi.

\subsubsection{Tổ chức kiểm thử}
Bộ kiểm thử API dùng Node.js test runner và chạy ứng dụng trên cổng động
trong cùng tiến trình. Các test tạo tài khoản có vai trò, gửi HTTP request
đến route và truy vấn PostgreSQL để xác minh dữ liệu. Phần lớn kiểm thử
được tích hợp với database đã cấu hình trong môi trường này; do đó, người
chạy lại cần database khả dụng và phải bảo đảm test không chạy nhầm vào
dữ liệu production. Tài liệu của Node.js mô tả API test runner và cách
khai báo test bất đồng bộ \cite{NodeTestDocs}.

\subsubsection{Kiểm soát an toàn dữ liệu}
Hệ thống chứa thông tin định danh tài khoản, câu trả lời và lịch sử thi.
Nguyên tắc cần áp dụng gồm chỉ truy xuất dữ liệu theo quyền sở hữu, tránh
ghi token/khóa vào log, mã hoá đường truyền khi triển khai, sao lưu có kiểm
soát và quy định thời hạn lưu dữ liệu. Khi gửi tài liệu cho nhà cung cấp
LLM, cần thông báo rõ đối tượng dữ liệu, mục đích và cơ chế lưu/huỷ. Các
nguyên tắc này phù hợp với khung quản trị AI và an ninh ứng dụng
\cite{NISTAIRMF,OWASPTop10,OWASPAPI,OWASPLLM}; chúng chưa được xem là đã xác
minh nếu chưa có kiểm toán chuyên biệt.

% ---------------------------------------------------------------------
%  CHƯƠNG 3 — KẾT QUẢ, ĐÁNH GIÁ VÀ KẾT LUẬN
% ---------------------------------------------------------------------
\chapter{Kết quả, đánh giá và kết luận}
\label{ch:ketqua}

\begin{muctieu}
Chương này đối chiếu mục tiêu với sản phẩm đã triển khai, trình bày bằng
chứng kiểm thử, thảo luận giới hạn và tổng kết hướng phát triển.
\end{muctieu}

\section{Kết quả đạt được}

Các mục tiêu ở Chương~\ref{ch:modau} được đối chiếu tại
Bảng~\ref{tab:doichieu}. Sản phẩm hiện bao gồm ứng dụng web, API và lược đồ
cơ sở dữ liệu; các luồng tạo câu hỏi, lập đề, duyệt đề, thi trực tuyến và
lớp học có mã nguồn cùng kiểm thử. Thành phần tạo câu hỏi hỗ trợ gọi mô
hình ngôn ngữ tùy cấu hình và có nhánh dự phòng. Các thành phần này chứng
minh mức độ hoàn thiện chức năng của nguyên mẫu, không chứng minh chất lượng
chuyên môn của câu hỏi do máy tạo.

\begin{table}[H]
\centering
\caption{Đối chiếu mục tiêu và kết quả triển khai.}
\label{tab:doichieu}
\begin{tabular*}{\textwidth}{@{\extracolsep{\fill}}p{0.31\textwidth}p{0.53\textwidth}p{0.10\textwidth}}
\toprule
\rowcolor{gray100}\bfseries\color{daunavy}Mục tiêu & \bfseries\color{daunavy}Bằng chứng trong dự án & \bfseries\color{daumaroon}Trạng thái \\
\midrule
Ứng dụng web nhiều thành phần & React/TypeScript, Express API, PostgreSQL và Docker Compose. & Đạt \\
Hỗ trợ sinh câu hỏi từ học liệu & Trích xuất PDF/DOCX, tích hợp OpenAI tùy chọn, sinh dự phòng và kiểm tra đầu ra. & Đạt một phần \\
Quản lý ngân hàng và đề & CRUD câu hỏi/đề, chọn theo loại và độ khó, tạo mã đề, quy trình duyệt. & Đạt \\
Tổ chức thi và lớp học & Mã phòng/lớp, lưu lượt làm bài, giám sát, tiếp tục bài và tự chấm trắc nghiệm. & Đạt \\
Kiểm thử chức năng và build & 35/35 bài kiểm thử API đạt; frontend build/type-check thành công. & Đạt \\
Đánh giá định lượng chất lượng AI & Chưa có tập chuẩn gán nhãn, phép đo chất lượng sinh hoặc benchmark OCR. & Chưa đạt \\
\bottomrule
\end{tabular*}
\end{table}

\section{Độ đo và bằng chứng kiểm thử}

Vì đồ án không huấn luyện mô hình và chưa có tập kiểm thử chuẩn cho câu hỏi
hay OCR, không thể báo cáo \en{accuracy}, precision, recall, F1 hoặc độ
chính xác OCR như kết quả thực nghiệm. Bằng chứng định lượng hiện có là số
bài kiểm thử phần mềm và kết quả build được trình bày trong
Bảng~\ref{tab:kiemthu}.

\begin{table}[H]
\centering
\caption{Kết quả xác minh phần mềm trong lần kiểm tra báo cáo.}
\label{tab:kiemthu}
\begin{tabular*}{\textwidth}{@{\extracolsep{\fill}}p{0.36\textwidth}p{0.36\textwidth}p{0.20\textwidth}}
\toprule
\rowcolor{gray100}\bfseries\color{daunavy}Hạng mục & \bfseries\color{daunavy}Kết quả & \bfseries\color{daumaroon}Đánh giá \\
\midrule
Kiểm thử API & 35 bài; 35 đạt, 0 lỗi, 0 bỏ qua. & Đạt \\
Build frontend & TypeScript và Vite build thành công. & Đạt \\
Sinh câu hỏi không có LLM & Có bài kiểm thử xác nhận nhánh sinh dự phòng trả ứng viên. & Đạt theo kiểm thử chức năng \\
Chất lượng nội dung AI & Chưa đối chiếu với bộ câu hỏi chuẩn hoặc đánh giá chuyên gia có thang điểm. & Chưa đo \\
Hiệu năng và tải & Chưa có benchmark thời gian phản hồi hoặc tải đồng thời. & Chưa đo \\
\bottomrule
\end{tabular*}
\end{table}

Backend trả về trường \code{confidence} theo nhánh xử lý, nhưng giá trị này
được gán trong mã nguồn và chưa được hiệu chuẩn bằng tập dữ liệu có nhãn.
Do đó, báo cáo không diễn giải trường này thành xác suất đúng hay chỉ số
đánh giá mô hình. Kiểm thử tự động tập trung vào cấu trúc đầu ra và nghiệp
vụ API; giảng viên vẫn cần kiểm tra nội dung câu hỏi, đáp án và bằng chứng
trước khi đưa vào đề thật.

\subsection{Ma trận kiểm thử chức năng}
Bảng~\ref{tab:ma-tran-kiem-thu} tóm tắt các nhóm hành vi được kiểm tra
trong bộ test hiện có. Mục đích là chỉ ra loại lỗi có thể được phát hiện,
không thay thế báo cáo từng lần chạy của CI. Tất cả 35 test hiện hữu đều
được thực thi thành công trong lần xác minh. Các yêu cầu không xuất hiện
trong test hiện hữu được giữ ở trạng thái chưa xác minh.
\begin{table}[H]
\centering
\caption{Ma trận kiểm thử chức năng theo rủi ro nghiệp vụ.}
\label{tab:ma-tran-kiem-thu}
\begin{tabularx}{\textwidth}{>{\bfseries\color{daunavy}}p{0.12\textwidth} X p{0.28\textwidth}}
\toprule
\rowcolor{gray100}Mã & Ca kiểm thử/điều kiện & Kết quả xác minh \\
\midrule
T01 & Health check và CORS preflight. & Đạt trong bộ test. \\
T02 & Cập nhật vai trò người dùng; token hiện hành và quyền được làm mới. & Đạt trong bộ test. \\
T03 & Sinh viên không truy cập phân tích quản trị; đăng ký công khai không tự cấp teacher/admin. & Đạt trong bộ test. \\
T04 & CRUD tài khoản quản trị; phân trang và bulk import. & Đạt trong bộ test. \\
T05 & Tạo, lưu và cập nhật đề; sinh mã đề với đáp án ánh xạ đúng. & Đạt trong bộ test. \\
T06 & Phân tích tài liệu có câu hỏi/đáp án; dữ liệu rỗng bị từ chối. & Đạt trong bộ test. \\
T07 & Không có LLM vẫn có ứng viên dự phòng; câu hỏi phải phù hợp văn bản nguồn. & Đạt theo ca hiện có. \\
T08 & Câu hỏi trắc nghiệm có lựa chọn riêng biệt, đáp án và căn cứ nguồn. & Đạt theo ca hiện có. \\
T09 & Phòng thi random đề, lưu bài, tự chấm trắc nghiệm và tiếp tục lượt. & Đạt trong bộ test. \\
T10 & Lớp học: tạo, tham gia bằng mã, danh sách thành viên và kiểm soát quyền. & Đạt trong bộ test. \\
T11 & OCR tiếng Việt, kiểm thử tải, tương phản/khả năng tiếp cận và trình duyệt thiết bị thật. & Chưa có bằng chứng. \\
\bottomrule
\end{tabularx}
\end{table}

\subsection{Diễn giải các loại kết quả}
Kiểm thử API kiểm chứng mã phản hồi, dữ liệu trả về và một phần trạng thái
PostgreSQL. Kiểm thử tích hợp chạm vào nhiều thành phần cùng lúc, nên có
thể phát hiện lỗi phối hợp mà test đơn vị của một hàm không phát hiện.
Tuy vậy, các bài test được viết cho những kịch bản xác định, không bao phủ
tất cả đầu vào, trình duyệt, nhà cung cấp OAuth hay cấu hình vận hành.

Lệnh build frontend xác nhận TypeScript không báo lỗi trong phạm vi
project build và bundler đóng gói được các tài nguyên. Nó không đo khả
năng sử dụng hoặc độ chính xác giao diện. Không có lượt test người dùng
được tiến hành, do vậy báo cáo không tự đưa ra kết luận về mức độ hài lòng
hoặc giảm thời gian biên soạn của giảng viên.

\subsection{Thiết kế thực nghiệm đánh giá AI đề xuất}
Để đánh giá chất lượng sinh câu hỏi trong tương lai, nên lập tập kiểm thử
đại diện cho nhiều môn và định dạng tài liệu, sau đó giảng viên độc lập
gán nhãn đúng/sai cho câu hỏi, đáp án, bằng chứng, độ rõ ràng, mức độ phù
hợp và độ khó. Một phần mẫu nên được hai người chấm để tính mức nhất quán
liên giám khảo; bất đồng cần có quy trình phân xử. Cần tách tài liệu đã
dùng để chỉnh prompt khỏi tập đánh giá cuối cùng nhằm hạn chế rò rỉ thông
tin vào phép đo.

Các câu hỏi trắc nghiệm có thể báo cáo tỷ lệ đáp án đúng được chuyên gia
xác nhận, tỷ lệ bằng chứng hỗ trợ phát biểu, tỷ lệ lựa chọn hợp lệ và
độ trùng lặp. Kết quả OCR có thể đo tỷ lệ lỗi ký tự và từ trên ảnh được
gán nhãn; điểm số OCR không thay cho đánh giá hậu quả của một lỗi OCR lên
câu hỏi sinh ra. Đánh giá phải phân tầng theo ngôn ngữ, loại tài liệu và
chất lượng scan, đồng thời công bố số mẫu và khoảng tin cậy.

\begin{table}[H]
\centering
\caption{Đề cương đo lường dành cho một đánh giá AI độc lập trong tương lai.}
\label{tab:ai-eval-plan}
\begin{tabularx}{\textwidth}{>{\bfseries\color{daunavy}}p{0.21\textwidth} X X}
\toprule
\rowcolor{gray100}Khía cạnh & Cách đo đề xuất & Đối tượng đánh giá \\
\midrule
Độ đúng nội dung & Tỷ lệ phát biểu/đáp án được chuyên gia xác nhận trên tập chưa dùng để phát triển. & Mỗi câu sinh và đáp án. \\
Bám nguồn & Tỷ lệ câu có bằng chứng nguyên văn, định vị đúng và đủ để hỗ trợ đáp án. & Cặp câu hỏi -- nguồn. \\
Chất lượng sư phạm & Thang điểm rõ ràng, một đáp án đúng, distractor hợp lý, ngôn ngữ và mức độ nhận thức. & Giảng viên chấm độc lập. \\
Độ đa dạng & Tỷ lệ trùng lặp nội dung và phân bố ý/chủ đề trong tập câu hỏi. & Tập đề sinh từ nhiều tài liệu. \\
OCR & Character/Word Error Rate trên ảnh quét có transcript chuẩn. & Từng nhóm ngôn ngữ và chất lượng ảnh. \\
Độ trễ/chi phí & Thời gian p50/p95, lỗi/time-out và chi phí theo tài liệu/câu hỏi. & Pipeline và từng nhà cung cấp. \\
\bottomrule
\end{tabularx}
\end{table}

Mỗi độ đo cần được báo cáo theo phương pháp đã định trước; không gộp nhiều
nhà cung cấp hoặc nhiều loại tài liệu thành một con số duy nhất nếu phân bố
mẫu khác nhau. Trường \code{confidence} hiện được gán theo nhánh xử lý và
chưa được hiệu chuẩn, nên không thể thay thế xác nhận chuyên gia hoặc
khoảng tin cậy thống kê.

\section{Bàn luận}

Điểm mạnh của thiết kế là có thể chạy quy trình nghiệp vụ liên tục từ tài
liệu và ngân hàng câu hỏi tới đề thi, lượt thi và kết quả; nhánh dự phòng
cũng cho phép kiểm thử luồng sinh mà không phụ thuộc khóa dịch vụ ngoài.
Việc lưu bằng chứng nguồn và hỗ trợ giảng viên rà soát là phù hợp với yêu
cầu truy xuất nội dung, nhưng không loại bỏ hoàn toàn nguy cơ câu hỏi sai
hoặc trích xuất sai.

Kết quả 35/35 bài đạt phản ánh phạm vi các ca kiểm thử đã viết. Nó không
đồng nghĩa hệ thống không còn lỗi, cũng không đo độ chính xác nội dung do
AI tạo. Build frontend thành công xác nhận khả năng biên dịch tại môi trường
đã kiểm tra, nhưng chưa xác nhận triển khai lên máy chủ thực tế.

\subsection{Đánh giá theo các vai trò}
Đối với giảng viên, việc gộp ngân hàng, bộ lọc ma trận và lịch sử đề làm
nền tảng cho khả năng tái sử dụng nội dung. Giá trị thực tiễn sẽ phụ thuộc
vào độ chính xác của căn cứ và thời gian phải chỉnh sửa câu AI; đồ án chưa
thực hiện đo trước/sau để kết luận tiết kiệm thời gian. Đối với sinh viên,
lưu vị trí và cho tiếp tục bài làm hỗ trợ tính liên tục, còn mã truy cập
tách việc dự thi khỏi việc gửi liên kết trực tiếp đến từng tài khoản.
Đối với quản trị viên, quản lý vai trò tập trung thuận tiện nhưng làm tăng
hậu quả nếu tài khoản có quyền cao bị chiếm dụng.

\subsection{Đánh đổi kiến trúc}
Tích hợp logic AI trong backend giúp sử dụng cơ sở dữ liệu và phân quyền
đã có, giảm số dịch vụ phải vận hành. Mặt trái là tác vụ trích xuất/gọi
LLM tốn thời gian có thể giữ request HTTP lâu và dùng bộ nhớ khi file được
đọc trong memory. Việc tách AI thành worker hoặc hàng đợi có thể cải thiện
độ bền và khả năng quan sát ở giai đoạn sau, nhưng đòi hỏi cơ chế trạng
thái job, retry an toàn, idempotency và giới hạn chi phí; chưa nên xem là
thành phần đã hiện thực.

\section{Hạn chế}

\begin{itemize}[leftmargin=1.6em, itemsep=3pt]
  \item Chưa có bộ dữ liệu chuẩn để đo độ đúng, độ bao phủ và độ đa dạng của
        câu hỏi sinh tự động.
  \item Chất lượng đầu ra phụ thuộc mô hình/API ngoài khi được cấu hình; bộ
        sinh dự phòng dựa trên quy tắc có thể tạo câu hỏi kém tự nhiên hoặc
        chưa đủ sâu.
  \item OCR chịu ảnh hưởng của chất lượng scan và bố cục; nhánh hiện tại
        dùng dữ liệu nhận dạng tiếng Anh, chưa có benchmark cho tài liệu
        tiếng Việt hoặc tập ảnh được gán nhãn.
  \item Câu tự luận cần được giảng viên chấm; giao diện có trường rubric và
        phản hồi nhưng báo cáo chưa xác nhận một bộ chấm tự động bằng AI.
  \item Chưa thực hiện kiểm thử tải, kiểm thử người dùng, đánh giá bảo mật
        triển khai và đo khả năng mở rộng.
  \item Các thư mục \code{ai-services/} và \code{ml-models/} mới là khung dự
        kiến; OCR và sinh câu hỏi hiện được gọi từ backend API.
\end{itemize}

% ---------------------------------------------------------------------
%  NỘI DUNG CHƯƠNG 3 — KẾT LUẬN VÀ HƯỚNG PHÁT TRIỂN
% ---------------------------------------------------------------------
\section{Kết luận và hướng phát triển}
\label{sec:ketluan}

\begin{muctieu}
Phần này tổng kết phạm vi hoàn thành, bài học rút ra và các hướng phát
triển có thể kiểm chứng cho AI ExamHub DAU.
\end{muctieu}

\subsection{Kết luận}

Đồ án đã khảo sát và xây dựng nguyên mẫu AI ExamHub DAU gồm frontend React,
API Express và cơ sở dữ liệu PostgreSQL. Hệ thống hỗ trợ quản lý câu hỏi,
tạo đề với trợ giúp từ nội dung tài liệu, duyệt đề, tổ chức phòng thi và
lớp học trực tuyến, lưu tiến độ làm bài và chấm câu trắc nghiệm. Bộ kiểm
thử API đạt 35/35 ca và frontend build thành công tại môi trường kiểm tra.
Tuy vậy, chưa có benchmark độc lập để kết luận về độ chính xác của sinh câu
hỏi, OCR hay hiệu năng; các đầu ra AI vẫn cần con người rà soát.

\subsection{Bài học rút ra}

Một thành phần AI hữu ích trong sản phẩm cần nhiều hơn một lời gọi mô hình:
cần kiểm soát nguồn vào, chuẩn hoá cấu trúc, kiểm tra đầu ra và cung cấp
cách để người dùng xác minh. Trong hệ thống tạo đề, việc lưu bằng chứng từ
tài liệu và giữ quyền duyệt cuối cùng cho giảng viên là các bước quan trọng
để hạn chế việc dùng nội dung không phù hợp.

Đồ án cũng cho thấy phải tách biệt kiểm thử chức năng với đánh giá mô hình.
Một endpoint chạy đúng không chứng minh câu hỏi hay hoặc đáp án chính xác.
Nếu tiếp tục thực hiện, cần lập bộ câu hỏi chuẩn, thống nhất tiêu chí chấm
với giảng viên và công bố phương pháp đo trước khi cải thiện prompt hay bộ
sinh dự phòng.

\subsection{Hướng phát triển}

\begin{itemize}[leftmargin=1.6em, itemsep=3pt]
  \item Xây dựng tập đánh giá có chuyên gia rà soát, đo mức bám nguồn, tính
        đúng của đáp án, độ rõ ràng, độ khó và mức trùng lặp.
  \item Hoàn thiện bộ xử lý OCR theo loại tài liệu, ngôn ngữ và chất lượng
        ảnh; lưu lại nguồn trích xuất để người dùng kiểm chứng.
  \item Thêm cơ chế duyệt/sửa câu hỏi AI, lịch sử phiên bản và chỉ đưa câu
        đã được phê duyệt vào ngân hàng chính thức.
  \item Đánh giá an toàn và quyền riêng tư khi gửi học liệu đến nhà cung cấp
        mô hình bên ngoài; cho phép cấu hình chính sách lưu trữ dữ liệu.
  \item Bổ sung kiểm thử tải, quan sát hệ thống và triển khai secrets an toàn;
        thay các giá trị phát triển bằng cấu hình vận hành phù hợp.
  \item Chỉ xem xét chấm tự luận có AI sau khi xác định rubric, bộ dữ liệu
        đối chiếu và quy trình giảng viên xác nhận kết quả.
\end{itemize}


% ---------------------------------------------------------------------
%  TÀI LIỆU THAM KHẢO VÀ PHỤ LỤC
% ---------------------------------------------------------------------
\backmatter

% \phantomsection tạo mốc neo riêng cho mục lục, nếu thiếu thì hyperref
% báo trùng neo giữa dấu trang và dấu trang cha.
\cleardoublepage
\phantomsection
\addcontentsline{toc}{chapter}{Tài liệu tham khảo}
\bibliographystyle{plain}
\bibliography{refs}

% ---------------------------------------------------------------------
%  PHỤ LỤC
% ---------------------------------------------------------------------
\chapter{Phụ lục}
\label{ch:phuluc}
\markboth{Phụ lục}{}

\section{Hướng dẫn cài đặt và chạy hệ thống}
\label{sec:huongdan-run}

Mã nguồn của hệ thống nằm trong thư mục \code{E:\textbackslash
AI-EXAMHUB-DAU}. Cấu hình thật được đọc từ các tệp môi trường không nên đưa
lên kho công khai. Để chạy stack phát triển bằng Docker Compose, tạo bản
cấu hình cục bộ từ ví dụ và khởi động các dịch vụ:

\begin{lstlisting}[language=bash, caption={Các lệnh khởi chạy stack phát triển.}]
cd E:\AI-EXAMHUB-DAU
docker compose up --build
\end{lstlisting}

Theo tệp Compose hiện có, frontend được công bố tại
\code{http://localhost:8080}, backend tại \code{http://localhost:3000/api}
và PostgreSQL tại cổng 5432. Docker Compose cho phép mô tả và khởi chạy
nhiều container từ một cấu hình \cite{DockerComposeDocs}. Các giá trị cổng
và mật khẩu development trong dự án chỉ phục vụ cài đặt cục bộ; phải đổi
trước khi triển khai có dữ liệu thật.

Để chạy từng phần, vào \code{backend/api}, cài phụ thuộc theo
\code{package-lock.json}, chuẩn bị PostgreSQL và cấu hình các biến môi
trường từ \code{.env.example}; sau đó chạy \code{npm run dev}. Frontend
được chạy trong thư mục \code{frontend/web} bằng \code{npm install} và
\code{npm run dev}; đặt \code{VITE\_API\_URL} trỏ tới backend đang chạy.
Các lệnh kiểm tra là \code{npm test} tại backend và
\code{npm run build} tại frontend. Nếu database schema chưa đầy đủ, cần
áp dụng các migration theo đúng thứ tự trong thư mục dự án. Không chạy
test bằng tài khoản hoặc database production.

\section{Tổng quan cơ sở dữ liệu}
\label{sec:db-overview}

Schema được khởi tạo từ \code{database/schema.sql} và bổ sung theo các
migration. Định danh chính dùng UUID; các nội dung linh hoạt như lựa chọn,
câu trả lời, danh sách ID hoặc metadata được lưu dạng JSONB; thời điểm
nghiệp vụ dùng \code{TIMESTAMPTZ}. Các khoá ngoại quy định vòng đời phụ
thuộc của bài nộp, kết quả, lượt thi và thành viên. Các quan hệ lõi được
minh hoạ ở Hình~\ref{fig:erd}; Bảng~\ref{tab:dulieu} trong
Chương~\ref{ch:thietke} là bản đồ thực thể ở mức khái quát.

\subsection{Quy ước đọc từ điển dữ liệu}
Trong các bảng dưới đây, PK là khoá chính; FK là khoá ngoại; UQ là giá trị
duy nhất; NN là bắt buộc; CK là ràng buộc miền/điều kiện. Kiểu JSONB hỗ trợ
dữ liệu JSON, nhưng không tự thay thế kiểm tra tính hợp lệ của cấu trúc
payload trong API. Một số cột được bổ sung qua migration hoặc
\code{ensureExamSchema()} lúc ứng dụng khởi động; vì vậy schema cuối cùng
đang chạy phải được đối chiếu với toàn bộ migration, không chỉ một file SQL.

\section{Từ điển dữ liệu}
\label{sec:data-dictionary}

\subsection{Bảng users -- tài khoản và vai trò}
Bảng \code{users} là gốc định danh của người dùng, giữ email đăng nhập,
thông tin hồ sơ và phạm vi quyền. Vai trò được giới hạn vào
\code{admin}, \code{teacher} hoặc \code{student}. Trường
\code{institutional\_id} hỗ trợ mã cán bộ/sinh viên và có chỉ mục duy nhất
riêng khi trường khác NULL. Không lưu mật khẩu rõ; API lưu giá trị băm ở
\code{password\_hash}.

\begin{table}[H]
\centering
\caption{Từ điển dữ liệu bảng users.}
\label{tab:dict-users}
\begin{tabularx}{\textwidth}{>{\ttfamily\color{daumaroon}}p{0.20\textwidth}p{0.17\textwidth}p{0.12\textwidth}X}
\toprule
\rowcolor{gray100}Trường & Kiểu & Ràng buộc & Ý nghĩa \\
\midrule
id & UUID & PK & Định danh tài khoản, mặc định sinh UUID. \\
email & VARCHAR(255) & NN, UQ & Địa chỉ đăng nhập, duy nhất. \\
password\_hash & TEXT & NN & Mật khẩu đã băm; không phải mật khẩu gốc. \\
full\_name & VARCHAR(255) & NN & Họ tên hiển thị. \\
role & VARCHAR(32) & NN, CK & Vai trò admin/teacher/student. \\
avatar\_url & TEXT & NULL & Đường dẫn hoặc nội dung ảnh đại diện theo API. \\
phone & VARCHAR(30) & NULL & Số điện thoại liên hệ. \\
institutional\_id & VARCHAR(64) & UQ, NULL & Mã định danh cơ sở đào tạo. \\
faculty & VARCHAR(255) & NULL & Khoa của người dùng. \\
department & VARCHAR(255) & NULL & Bộ môn hoặc đơn vị. \\
cohort & VARCHAR(64) & NULL & Khoá học. \\
class\_name & VARCHAR(128) & NULL & Lớp hành chính. \\
created\_at & TIMESTAMPTZ & NN & Thời điểm tạo tài khoản. \\
\bottomrule
\end{tabularx}
\end{table}

\subsection{Bảng questions -- ngân hàng câu hỏi}
Bảng \code{questions} giữ nội dung và metadata cần cho việc chọn câu vào đề.
Các trường liên quan loại câu, câu trả lời, mức nhận thức và nguồn có thể
được mở rộng qua migration. Lựa chọn và rubric có cấu trúc linh hoạt; nội
dung nguồn được lưu riêng ở \code{source\_evidence} để phục vụ rà soát và
truy nguyên. Câu hỏi có thể tồn tại độc lập với người tạo nếu tài khoản
gốc bị xoá.

\begin{table}[H]
\centering
\caption{Từ điển dữ liệu bảng questions.}
\label{tab:dict-questions}
\begin{tabularx}{\textwidth}{>{\ttfamily\color{daumaroon}}p{0.21\textwidth}p{0.17\textwidth}p{0.12\textwidth}X}
\toprule
\rowcolor{gray100}Trường & Kiểu & Ràng buộc & Ý nghĩa \\
\midrule
id & UUID & PK & Định danh câu hỏi. \\
subject & VARCHAR(255) & NN & Môn/chủ đề học thuật. \\
content & TEXT & NN & Lời dẫn hoặc nội dung câu hỏi. \\
choices & JSONB & NN & Danh sách lựa chọn; mặc định mảng rỗng. \\
answer\_index & INTEGER & NULL & Chỉ số lựa chọn đúng cho câu trắc nghiệm. \\
answer\_text & TEXT & NULL & Đáp án tham khảo dạng văn bản. \\
source\_evidence & TEXT & NULL & Trích đoạn nguồn làm căn cứ cho câu hỏi. \\
cognitive\_level & VARCHAR(32) & NULL & Mức nhận thức như understanding/application. \\
rubric & JSONB & NULL & Cấu trúc tiêu chí chấm tự luận nếu có. \\
difficulty & VARCHAR(16) & NN, CK & easy, medium hoặc hard; mặc định medium. \\
type & VARCHAR(32) & Schema mở rộng & Loại multipleChoice/essay; được bổ sung trong schema runtime/migration. \\
created\_by & UUID & FK, NULL & Tài khoản tạo; xoá tài khoản đặt NULL. \\
created\_at & TIMESTAMPTZ & NN & Thời điểm tạo câu hỏi. \\
\bottomrule
\end{tabularx}
\end{table}

\subsection{Bảng exams -- đề thi}
Bảng \code{exams} lưu thông tin định danh và trạng thái vòng đời của đề.
Các chi tiết cấu hình đề được lưu trong \code{metadata} dạng JSONB; các
liên kết câu hỏi có thứ tự được chuẩn hoá ở bảng \code{exam\_questions}.
Schema cũng có \code{question\_ids} nhằm tương thích với các luồng cũ.
Khi bảo trì cần xác định một nguồn dữ liệu chuẩn để tránh hai biểu diễn
bị lệch nhau.

\begin{table}[H]
\centering
\caption{Từ điển dữ liệu bảng exams.}
\label{tab:dict-exams}
\begin{tabularx}{\textwidth}{>{\ttfamily\color{daumaroon}}p{0.21\textwidth}p{0.17\textwidth}p{0.12\textwidth}X}
\toprule
\rowcolor{gray100}Trường & Kiểu & Ràng buộc & Ý nghĩa \\
\midrule
id & UUID & PK & Định danh đề. \\
title & VARCHAR(255) & NN & Tên đề thi. \\
subject & VARCHAR(255) & NN & Môn học. \\
status & VARCHAR(16) & NN, CK & Trạng thái phát hành ở schema. \\
workflow\_status & VARCHAR(16) & NN/default & Trạng thái quy trình: draft/review/approved/published trong ứng dụng. \\
question\_ids & JSONB & NN & Danh sách ID tương thích với luồng legacy. \\
metadata & JSONB & NN & Thời lượng, tổng điểm, ma trận, năm học, rubric và thông tin đề. \\
created\_by & UUID & FK, NULL & Tài khoản tạo; xoá người tạo đặt NULL. \\
created\_at & TIMESTAMPTZ & NN & Thời điểm tạo. \\
updated\_at & TIMESTAMPTZ & NN & Thời điểm cập nhật gần nhất. \\
\bottomrule
\end{tabularx}
\end{table}

\subsection{Bảng exam\_questions -- câu hỏi thuộc đề}
Bảng nối biểu diễn quan hệ nhiều--nhiều giữa đề và câu hỏi, đồng thời
lưu thứ tự hiển thị. Một câu hỏi có thể tái sử dụng trong nhiều đề; xoá đề
sẽ cascade xoá liên kết nhưng không xoá câu hỏi.
\begin{table}[H]
\centering
\caption{Từ điển dữ liệu bảng exam\_questions.}
\label{tab:dict-exam-questions}
\begin{tabularx}{\textwidth}{>{\ttfamily\color{daumaroon}}p{0.21\textwidth}p{0.17\textwidth}p{0.12\textwidth}X}
\toprule
\rowcolor{gray100}Trường & Kiểu & Ràng buộc & Ý nghĩa \\
\midrule
exam\_id & UUID & PK, FK & Đề chứa câu hỏi; xoá đề cascade liên kết. \\
question\_id & UUID & PK, FK & Câu hỏi; không cho xoá nếu vẫn thuộc đề. \\
position & INTEGER & NN & Vị trí câu trong đề. \\
\bottomrule
\end{tabularx}
\end{table}

\subsection{Bảng exam\_rooms -- phòng thi}
Phòng thi gắn với một đề chính và có thể gắn nhiều đề bổ sung qua bảng
\code{exam\_room\_exams}. Mã tham gia được đảm bảo duy nhất. Thời điểm mở,
đóng và trạng thái điều khiển khả năng vào phòng.
\begin{table}[H]
\centering
\caption{Từ điển dữ liệu bảng exam\_rooms.}
\label{tab:dict-exam-rooms}
\begin{tabularx}{\textwidth}{>{\ttfamily\color{daumaroon}}p{0.21\textwidth}p{0.17\textwidth}p{0.12\textwidth}X}
\toprule
\rowcolor{gray100}Trường & Kiểu & Ràng buộc & Ý nghĩa \\
\midrule
id & UUID & PK & Định danh phòng thi. \\
exam\_id & UUID & FK, NN & Đề chính của phòng; xoá đề cascade phòng. \\
name & VARCHAR(255) & NN & Tên phòng. \\
starts\_at & TIMESTAMPTZ & NULL & Thời điểm dự kiến bắt đầu. \\
ends\_at & TIMESTAMPTZ & NULL & Thời điểm kết thúc/đóng nhận bài. \\
status & VARCHAR(16) & NN, CK & scheduled, active hoặc closed. \\
access\_code & VARCHAR(8) & NN, UQ & Mã tham gia ngẫu nhiên duy nhất. \\
created\_by & UUID & FK, NULL & Tài khoản tạo phòng. \\
\bottomrule
\end{tabularx}
\end{table}

\subsection{Bảng exam\_room\_exams -- đề có thể phân cho phòng}
Bảng liên kết phòng với một hoặc nhiều đề được giảng viên chọn. Mỗi đề có
vị trí trong danh sách; một lượt thi chọn đề được gán cho sinh viên và
lưu ID đề ở \code{exam\_attempts}.
\begin{table}[H]
\centering
\caption{Từ điển dữ liệu bảng exam\_room\_exams.}
\label{tab:dict-room-exams}
\begin{tabularx}{\textwidth}{>{\ttfamily\color{daumaroon}}p{0.21\textwidth}p{0.17\textwidth}p{0.12\textwidth}X}
\toprule
\rowcolor{gray100}Trường & Kiểu & Ràng buộc & Ý nghĩa \\
\midrule
room\_id & UUID & PK, FK & Phòng thi; xoá phòng cascade liên kết. \\
exam\_id & UUID & PK, FK & Đề được đưa vào danh sách phân bổ. \\
position & INTEGER & NN, UQ theo phòng & Vị trí đề; duy nhất trong một phòng. \\
\bottomrule
\end{tabularx}
\end{table}

\subsection{Bảng exam\_attempts -- lượt dự thi}
Lượt thi lưu trạng thái làm bài độc lập với định nghĩa đề. Bộ câu hỏi được
giao và đáp án lưu ở JSONB để tiếp tục bài làm; khoá duy nhất
\code{(room\_id, student\_id)} bảo đảm một lượt chính cho sinh viên trong
phòng theo schema hiện tại.
\begin{table}[H]
\centering
\caption{Từ điển dữ liệu bảng exam\_attempts.}
\label{tab:dict-attempts}
\begin{tabularx}{\textwidth}{>{\ttfamily\color{daumaroon}}p{0.21\textwidth}p{0.17\textwidth}p{0.12\textwidth}X}
\toprule
\rowcolor{gray100}Trường & Kiểu & Ràng buộc & Ý nghĩa \\
\midrule
id & UUID & PK & Định danh lượt thi. \\
room\_id & UUID & FK, NN & Phòng thi chứa lượt. \\
exam\_id & UUID & FK, NN & Đề được giao. \\
student\_id & UUID & FK, NN & Sinh viên; xoá tài khoản cascade lượt. \\
question\_set & JSONB & NN & Snapshot câu hỏi an toàn của lượt thi. \\
answers & JSONB & NN & Câu trả lời theo ID; mặc định object rỗng. \\
marked\_questions & JSONB & NN & Danh sách câu được đánh dấu xem lại. \\
current\_question & INTEGER & NN & Vị trí hiện tại để tiếp tục bài. \\
started\_at & TIMESTAMPTZ & NN & Thời điểm bắt đầu. \\
deadline\_at & TIMESTAMPTZ & NN & Hạn hoàn thành bài. \\
submitted\_at & TIMESTAMPTZ & NULL & Thời điểm nộp bài. \\
submission\_id & UUID & FK, UQ, NULL & Bài nộp đã tạo cho lượt thi. \\
updated\_at & TIMESTAMPTZ & NN & Thời điểm cập nhật gần nhất. \\
\bottomrule
\end{tabularx}
\end{table}

\subsection{Bảng exam\_attempt\_events -- sự kiện lượt thi}
Bảng sự kiện hỗ trợ giám sát một số dấu hiệu thao tác trình duyệt. Schema
hiện tại định nghĩa loại sự kiện \code{tab\_hidden}; đây không phải bằng
chứng đầy đủ về hành vi của sinh viên, cũng không thay thế cơ chế giám sát
trực tiếp hoặc quy chế thi.
\begin{table}[H]
\centering
\caption{Từ điển dữ liệu bảng exam\_attempt\_events.}
\label{tab:dict-attempt-events}
\begin{tabularx}{\textwidth}{>{\ttfamily\color{daumaroon}}p{0.21\textwidth}p{0.17\textwidth}p{0.12\textwidth}X}
\toprule
\rowcolor{gray100}Trường & Kiểu & Ràng buộc & Ý nghĩa \\
\midrule
id & UUID & PK & Định danh sự kiện. \\
attempt\_id & UUID & FK, NN & Lượt thi; xoá lượt cascade sự kiện. \\
event\_type & VARCHAR(32) & NN, CK & Loại sự kiện cho phép, hiện gồm tab\_hidden. \\
occurred\_at & TIMESTAMPTZ & NN & Thời điểm ghi nhận sự kiện. \\
\bottomrule
\end{tabularx}
\end{table}

\subsection{Bảng submissions -- bài nộp}
Submission là bản ghi nộp bài có cấu trúc thống nhất cho lịch sử và chấm
điểm. Câu trả lời được giữ trong JSONB; trạng thái phân biệt bài đã nộp,
đang chấm hoặc đã chấm. Sinh viên có thể bị đặt NULL khi tài khoản bị
xoá theo chính sách schema.
\begin{table}[H]
\centering
\caption{Từ điển dữ liệu bảng submissions.}
\label{tab:dict-submissions}
\begin{tabularx}{\textwidth}{>{\ttfamily\color{daumaroon}}p{0.21\textwidth}p{0.17\textwidth}p{0.12\textwidth}X}
\toprule
\rowcolor{gray100}Trường & Kiểu & Ràng buộc & Ý nghĩa \\
\midrule
id & UUID & PK & Định danh bài nộp. \\
exam\_id & UUID & FK, NN & Đề được nộp. \\
student\_id & UUID & FK, NULL & Sinh viên nộp. \\
answers & JSONB & NN & Nội dung trả lời đã nộp. \\
status & VARCHAR(16) & NN, CK & submitted, grading hoặc graded. \\
submitted\_at & TIMESTAMPTZ & NN & Thời điểm nộp. \\
\bottomrule
\end{tabularx}
\end{table}

\subsection{Bảng results -- kết quả chấm}
Mỗi bài nộp có tối đa một bản ghi kết quả. Kết quả giữ điểm, phản hồi và
thời điểm chấm. Mức điểm tối đa được xác định từ metadata của đề và được
kiểm tra trong luồng chấm.
\begin{table}[H]
\centering
\caption{Từ điển dữ liệu bảng results.}
\label{tab:dict-results}
\begin{tabularx}{\textwidth}{>{\ttfamily\color{daumaroon}}p{0.21\textwidth}p{0.17\textwidth}p{0.12\textwidth}X}
\toprule
\rowcolor{gray100}Trường & Kiểu & Ràng buộc & Ý nghĩa \\
\midrule
id & UUID & PK & Định danh kết quả. \\
submission\_id & UUID & FK, NN, UQ & Bài nộp; một kết quả cho mỗi bài. \\
score & NUMERIC(5,2) & NN & Điểm số, mặc định 0. \\
feedback & TEXT & NULL & Nhận xét của người chấm. \\
graded\_at & TIMESTAMPTZ & NN & Thời điểm ghi/cập nhật điểm. \\
\bottomrule
\end{tabularx}
\end{table}

\subsection{Bảng rubrics -- tiêu chí chấm}
Bảng \code{rubrics} gắn bộ tiêu chí với đề thi. Tiêu chí có dạng JSONB để
mô tả cấu trúc rubric mà không tạo bảng con cho từng cấp điểm. Việc lưu
rubric không đồng nghĩa đã có công cụ AI tự động áp dụng rubric.
\begin{table}[H]
\centering
\caption{Từ điển dữ liệu bảng rubrics.}
\label{tab:dict-rubrics}
\begin{tabularx}{\textwidth}{>{\ttfamily\color{daumaroon}}p{0.21\textwidth}p{0.17\textwidth}p{0.12\textwidth}X}
\toprule
\rowcolor{gray100}Trường & Kiểu & Ràng buộc & Ý nghĩa \\
\midrule
id & UUID & PK & Định danh bộ rubric. \\
exam\_id & UUID & FK, NN & Đề áp dụng; xoá đề cascade rubric. \\
title & VARCHAR(255) & NN & Tên tiêu chí chấm. \\
criteria & JSONB & NN & Danh sách tiêu chí/mức điểm; mặc định mảng rỗng. \\
\bottomrule
\end{tabularx}
\end{table}

\subsection{Bảng classroom\_rooms -- lớp học trực tuyến}
Lớp học lưu thông tin buổi học, chế độ kết nối và quyền nhận thành viên.
Mã tham gia là duy nhất; tài khoản tạo lớp là bắt buộc. Google Meet URL
được kiểm tra ở API khi chế độ kết nối là Meet.
\begin{table}[H]
\centering
\caption{Từ điển dữ liệu bảng classroom\_rooms.}
\label{tab:dict-classroom-rooms}
\begin{tabularx}{\textwidth}{>{\ttfamily\color{daumaroon}}p{0.21\textwidth}p{0.17\textwidth}p{0.12\textwidth}X}
\toprule
\rowcolor{gray100}Trường & Kiểu & Ràng buộc & Ý nghĩa \\
\midrule
id & UUID & PK & Định danh lớp. \\
name & VARCHAR(160) & NN & Tên lớp học. \\
description & TEXT & NN & Mô tả, mặc định chuỗi rỗng. \\
mode & VARCHAR(20) & NN, CK & google\_meet hoặc direct. \\
meet\_url & TEXT & NULL & Liên kết Google Meet nếu chế độ tương ứng. \\
access\_code & VARCHAR(8) & NN, UQ & Mã tham gia lớp duy nhất. \\
status & VARCHAR(16) & NN, CK & scheduled, live hoặc ended. \\
admissions\_open & BOOLEAN & NN & Có nhận sinh viên mới hay không. \\
started\_at & TIMESTAMPTZ & NULL & Thời điểm bắt đầu buổi học. \\
ended\_at & TIMESTAMPTZ & NULL & Thời điểm kết thúc buổi học. \\
created\_by & UUID & FK, NN & Giảng viên/quản trị viên tạo lớp. \\
created\_at & TIMESTAMPTZ & NN & Thời điểm tạo. \\
\bottomrule
\end{tabularx}
\end{table}

\subsection{Bảng classroom\_members -- thành viên lớp}
Bảng nối thành viên có khoá ghép để cùng một sinh viên không gia nhập
trùng một lớp. Lịch sử gia nhập được lưu theo thời gian.
\begin{table}[H]
\centering
\caption{Từ điển dữ liệu bảng classroom\_members.}
\label{tab:dict-classroom-members}
\begin{tabularx}{\textwidth}{>{\ttfamily\color{daumaroon}}p{0.21\textwidth}p{0.17\textwidth}p{0.12\textwidth}X}
\toprule
\rowcolor{gray100}Trường & Kiểu & Ràng buộc & Ý nghĩa \\
\midrule
classroom\_id & UUID & PK, FK & Lớp học; xoá lớp cascade thành viên. \\
student\_id & UUID & PK, FK & Tài khoản sinh viên; xoá tài khoản cascade quan hệ. \\
joined\_at & TIMESTAMPTZ & NN & Thời điểm tham gia. \\
\bottomrule
\end{tabularx}
\end{table}

\section{Danh mục API theo nhóm nghiệp vụ}
\label{sec:api-catalog}

Các endpoint dưới đây được tổng hợp từ route và tài liệu API trong dự án.
API có tiền tố \code{/api}; các endpoint đánh dấu “cần xác thực” yêu cầu
JWT, còn thao tác quản trị/giảng viên có kiểm tra vai trò. Hợp đồng HTTP
có thể được mô tả bằng OpenAPI để tạo tài liệu và kiểm thử nhất quán
\cite{OpenAPISpec}.

\subsection{Xác thực, tài khoản và ngân hàng câu hỏi}
\begin{table}[H]
\centering
\caption{Endpoint xác thực, người dùng và câu hỏi.}
\label{tab:api-auth-questions}
\begin{tabularx}{\textwidth}{>{\ttfamily\color{daumaroon}}p{0.13\textwidth}p{0.39\textwidth}X}
\toprule
\rowcolor{gray100}Method & Endpoint & Mục đích / quyền \\
\midrule
POST & \code{/auth/register} & Đăng ký tài khoản sinh viên. \\
POST & \code{/auth/login} & Đăng nhập bằng email/mật khẩu. \\
GET & \code{/auth/me} & Lấy hồ sơ/token mới; cần xác thực. \\
PUT & \code{/auth/profile} & Cập nhật hồ sơ tài khoản; cần xác thực. \\
PUT & \code{/auth/password} & Đổi mật khẩu; cần xác thực. \\
GET & \code{/auth/google} & Khởi tạo đăng nhập Google nếu được cấu hình. \\
GET & \code{/auth/google/callback} & Nhận kết quả chuyển hướng OAuth. \\
GET & \code{/users} & Danh sách tài khoản/phân trang; admin. \\
POST & \code{/users} & Tạo tài khoản quản trị; admin. \\
GET & \code{/questions} & Tra cứu, lọc và phân trang câu hỏi. \\
POST/PUT/DELETE & \code{/questions} & Thêm/sửa/xoá câu hỏi; teacher/admin, xoá/sửa còn kiểm tra quyền sở hữu. \\
\bottomrule
\end{tabularx}
\end{table}

\subsection{Tài liệu và quản lý đề thi}
\begin{table}[H]
\centering
\caption{Endpoint xử lý tài liệu và vòng đời đề thi.}
\label{tab:api-exams}
\begin{tabularx}{\textwidth}{>{\ttfamily\color{daumaroon}}p{0.13\textwidth}p{0.42\textwidth}X}
\toprule
\rowcolor{gray100}Method & Endpoint & Mục đích / quyền \\
\midrule
POST & \code{/documents/upload} & Tải và trích xuất tài liệu; teacher/admin. \\
POST & \code{/documents/analyze} & Phân tích nội dung và trả ứng viên câu hỏi; teacher/admin. \\
POST & \code{/exams/generate} & Tạo đề từ câu hỏi đã chọn; teacher/admin. \\
GET & \code{/exams} & Danh sách đề theo phạm vi người dùng. \\
GET & \code{/exams/:id} & Chi tiết đề; câu trả về cho sinh viên ẩn đáp án nội bộ. \\
PUT & \code{/exams/:id} & Cập nhật metadata/trạng thái đề. \\
DELETE & \code{/exams/:id} & Xoá đề theo quyền. \\
POST & \code{/exams/:id/variants} & Tạo mã đề và ánh xạ lại đáp án. \\
POST & \code{/exams/:id/validate} & Kiểm tra cấu trúc/ma trận đề. \\
\bottomrule
\end{tabularx}
\end{table}

\subsection{Phòng thi, bài nộp và lớp học}
\begin{table}[H]
\centering
\caption{Endpoint tổ chức thi, chấm bài và quản lý lớp học.}
\label{tab:api-rooms-classrooms}
\begin{tabularx}{\textwidth}{>{\ttfamily\color{daumaroon}}p{0.13\textwidth}p{0.43\textwidth}X}
\toprule
\rowcolor{gray100}Method & Endpoint & Mục đích / quyền \\
\midrule
POST & \code{/exam-rooms} & Teacher/admin tạo phòng và mã tham gia. \\
GET & \code{/exam-rooms} & Teacher/admin liệt kê phòng được quản lý. \\
PATCH & \code{/exam-rooms/:id/status} & Mở/đóng phòng; kiểm tra chủ sở hữu. \\
GET & \code{/exam-rooms/:id/attempts} & Giám sát tiến độ và sự kiện; teacher/admin. \\
GET & \code{/exam-rooms/:id/grading} & Lấy các lượt thi đã nộp để chấm. \\
POST & \code{/exam-rooms/:code/start} & Sinh viên bắt đầu/tiếp tục bài bằng mã phòng. \\
PUT & \code{/exam-rooms/attempts/:id/answers} & Lưu câu trả lời và vị trí; chủ lượt thi. \\
POST & \code{/exam-rooms/attempts/:id/submit} & Nộp bài, tự chấm trắc nghiệm. \\
POST & \code{/exam-rooms/attempts/:id/grade} & Teacher/admin lưu điểm/phản hồi tự luận. \\
GET & \code{/exam-rooms/my-attempts} & Sinh viên xem lịch sử dự thi. \\
POST & \code{/classrooms} & Teacher/admin tạo lớp học. \\
POST & \code{/classrooms/join} & Sinh viên gia nhập bằng mã lớp. \\
GET/PATCH & \code{/classrooms/:id} & Xem lớp hoặc đổi trạng thái/nhận thành viên theo quyền. \\
\bottomrule
\end{tabularx}
\end{table}

\section{Danh mục migration dữ liệu}
\label{sec:migrations}

Các migration trong dự án ghi nhận việc mở rộng schema theo chức năng.
Bảng~\ref{tab:migrations} tóm tắt từ tên file; trước khi chạy trên dữ liệu
thật cần đọc nội dung migration, xác minh thứ tự, sao lưu và có phương án
rollback.
\begin{table}[H]
\centering
\caption{Danh mục migration hiện diện trong dự án.}
\label{tab:migrations}
\begin{tabularx}{\textwidth}{>{\ttfamily\color{daumaroon}}p{0.16\textwidth}p{0.47\textwidth}X}
\toprule
\rowcolor{gray100}Migration & Phạm vi thể hiện ở tên/tệp & Nhóm dữ liệu \\
\midrule
001 & \code{initial\_schema} & Schema ban đầu: người dùng, câu hỏi, đề và dữ liệu liên quan. \\
002 & \code{add\_metadata\_to\_exams} & Metadata cấu hình cho đề thi. \\
003 & \code{exam\_generation\_validation} & Trường phục vụ kiểm tra sinh đề. \\
004 & \code{add\_answer\_text\_to\_questions} & Đáp án dạng văn bản cho câu hỏi. \\
005 & \code{add\_question\_source\_evidence} & Bằng chứng trích từ tài liệu nguồn. \\
006 & \code{add\_essay\_cognitive\_rubric} & Mức độ nhận thức và rubric câu tự luận. \\
007 & \code{online\_exam\_rooms} & Phòng thi, lượt thi trực tuyến và sự kiện. \\
008 & \code{user\_profiles} & Trường hồ sơ tài khoản. \\
009 & \code{university\_user\_profiles} & Mã và thuộc tính hồ sơ đại học. \\
010 & \code{exam\_tab\_switch\_events} & Lịch sử sự kiện chuyển tab. \\
011 & \code{classrooms} & Lớp học, thành viên và mã tham gia. \\
012 & \code{classroom\_lifecycle} & Trạng thái vòng đời và nhận thành viên lớp. \\
\bottomrule
\end{tabularx}
\end{table}

\section{Ma trận kiểm thử bổ sung}
\label{sec:test-matrix}

Bên cạnh 35 ca tự động đã chạy, Bảng~\ref{tab:manual-tests} là bộ kịch bản
kiểm thử thủ công đề xuất cho buổi trình diễn. Đây không phải kết quả đã
chạy; cột trạng thái được để trống để người nộp cập nhật sau khi thực hiện.
\begin{table}[H]
\centering
\caption{Kịch bản kiểm thử thủ công trước khi bảo vệ.}
\label{tab:manual-tests}
\begin{tabularx}{\textwidth}{>{\bfseries\color{daumaroon}}p{0.10\textwidth}X X p{0.15\textwidth}}
\toprule
\rowcolor{gray100}Mã & Bước kiểm thử & Kết quả mong đợi & Trạng thái \\
\midrule
M01 & Đăng nhập lần lượt bằng vai trò sinh viên và giảng viên. & Menu và endpoint phù hợp vai trò; quyền không hợp lệ bị từ chối. & Chưa ghi nhận \\
M02 & Tải PDF có văn bản thật và kiểm tra phần text trích xuất. & Nội dung có thể xem trước; câu hỏi liên kết được với đoạn nguồn. & Chưa ghi nhận \\
M03 & Tải PDF quét chất lượng thấp hoặc file rỗng. & Có cảnh báo trích xuất; không sinh đề thành công giả. & Chưa ghi nhận \\
M04 & Tạo đề ma trận vượt số câu hiện có. & Hiện rõ số câu thiếu; không phát hành đề không hợp lệ. & Chưa ghi nhận \\
M05 & Gửi duyệt và chuyển trạng thái đề. & Trạng thái lưu sau khi tải lại trang. & Chưa ghi nhận \\
M06 & Vào phòng thi, trả lời một phần, tải lại rồi tiếp tục. & Đáp án/vị trí được khôi phục từ đúng lượt. & Chưa ghi nhận \\
M07 & Nộp bài có trắc nghiệm và tự luận. & Trắc nghiệm được chấm; tự luận chờ giảng viên chấm. & Chưa ghi nhận \\
M08 & Sinh viên thử xem API chi tiết đề/đáp án nội bộ. & API không trả đáp án đúng của bài thi. & Chưa ghi nhận \\
M09 & Gia nhập lớp đã kết thúc hoặc đóng nhận thành viên. & Từ chối người mới; thành viên đã có vẫn theo chính sách. & Chưa ghi nhận \\
M10 & Mở giao diện trên thiết bị di động, bàn phím và phóng đại. & Bố cục và thao tác chính còn sử dụng được. & Chưa ghi nhận \\
\bottomrule
\end{tabularx}
\end{table}

\section{Rủi ro, giới hạn và biện pháp giảm thiểu}
\label{sec:risk-register}

\begin{table}[H]
\centering
\caption{Sổ đăng ký rủi ro của nguyên mẫu AI ExamHub DAU.}
\label{tab:risk-register}
\begin{tabularx}{\textwidth}{>{\bfseries\color{daumaroon}}p{0.17\textwidth}X X}
\toprule
\rowcolor{gray100}Rủi ro & Tác động & Biện pháp đề xuất \\
\midrule
Đáp án AI sai & Sinh viên nhận nội dung sai; đề thi mất độ tin cậy. & Bắt buộc giảng viên xác minh câu, đáp án và nguồn trước phát hành. \\
Prompt injection từ tài liệu & Nội dung tải lên có thể chứa chỉ dẫn làm lệch yêu cầu cho mô hình. & Cô lập tài liệu khỏi instruction, loại hành động công cụ nhạy cảm, kiểm tra đầu ra và quyền duyệt của người. \\
Lộ dữ liệu qua API LLM & Học liệu/tài liệu có dữ liệu cá nhân đi ra nhà cung cấp. & Cảnh báo trước khi gửi, lọc dữ liệu nhạy cảm, cấu hình chính sách lưu giữ và nhà cung cấp được phép. \\
Lộ đáp án thi & Lộ đáp án làm mất công bằng kiểm tra. & Backend chỉ trả question view đã loại answer fields; test truy cập trực tiếp bằng vai trò sinh viên. \\
Chiếm đoạt tài khoản & Truy cập hồ sơ, đề, kết quả và dữ liệu lớp. & Secret mạnh, HTTPS, xác thực phù hợp, giới hạn quyền, nhật ký và quy trình đổi/thu hồi token. \\
Tệp độc hại/quá lớn & Cạn bộ nhớ, lỗi parser hoặc ảnh hưởng dịch vụ. & Giới hạn kích thước/định dạng, kiểm tra nội dung, cách ly xử lý và giới hạn tài nguyên. \\
OCR tiếng Việt kém & Nội dung nguồn bị nhận sai rồi lan truyền sang câu hỏi. & Định cấu hình ngôn ngữ phù hợp, hiển thị nguồn, đo lỗi OCR và cho phép người dùng hiệu chỉnh. \\
Database development & Mất dữ liệu hoặc lộ mật khẩu do cấu hình mặc định. & Thay mật khẩu, tách network, sao lưu/khôi phục kiểm chứng và dùng secret manager. \\
\bottomrule
\end{tabularx}
\end{table}

Sổ đăng ký là phân tích rủi ro thiết kế, không phải báo cáo kiểm thử xâm
nhập. Các rủi ro prompt injection và phân quyền dịch vụ API cần được đối
chiếu với khuyến nghị OWASP về ứng dụng web, API và LLM
\cite{OWASPTop10,OWASPAPI,OWASPLLM}; khung NIST giúp tổ chức việc nhận
diện, đo và xử lý rủi ro AI theo vòng đời \cite{NISTAIRMF}.

\section{Mục tự kiểm trước khi nộp}
\label{sec:checklist}

\begin{tukiem}
  \item Đã thay tên, MSSV, lớp, giảng viên hướng dẫn và năm học trên bìa.
  \item Đã rà soát mọi chương, mục lục, danh mục hình/bảng và số tham chiếu.
  \item Đã kiểm tra lại nội dung PDF/DOCX/OCR, các hạn chế kỹ thuật và phiên
        bản code được dùng để báo cáo.
  \item Đã chạy lại API tests, frontend build và ghi đúng phiên bản kết quả.
  \item Đã kiểm tra các đề xuất AI với tài liệu gốc; không coi confidence là
        xác suất đúng.
  \item Đã bổ sung ảnh chụp giao diện thật nếu giảng viên yêu cầu; các hình
        wireframe trong báo cáo được nhận diện là hình minh hoạ.
  \item Đã kiểm tra và ghi nguồn cho tối thiểu 15 tài liệu tham khảo có liên
        kết; mọi tài liệu trích dẫn đều có trong danh mục.
  \item Đã thay toàn bộ cấu hình development trước khi triển khai có dữ liệu
        thật.
  \item Đã xác nhận số trang sau khi biên dịch XeLaTeX; không coi số trang
        ước lượng trong source là số trang PDF.
  \item Đã xoá nội dung hướng dẫn và placeholder không thuộc bản nộp cuối.
\end{tukiem}

\section{Nhận xét của giảng viên hướng dẫn}
\label{sec:nhanxet}

\khungnhanxet[7.5cm]

\vspace{0.8em}
\noindent\begin{minipage}[t]{0.5\textwidth}
\color{daunavy}
{\color{daumaroon}\bfseries Điểm số}

\vspace{0.7em}
\dongchamnhan{Bằng số:}

\vspace{0.9em}
\dongchamnhan{Bằng chữ:}
\end{minipage}\hfill
\okyten{\dauGVHD}{Giảng viên hướng dẫn}


\end{document}
